% Generated mechanically from frozen input inputs/p06/ja/more-groupoids.tex.
% Do not edit this generated file; edit the bound locale source instead.

% OK, start here.
%

\setcounter{chapter}{39}
\chapter{亜群スキームの続論}



\phantomsection
\label{more-groupoids-section-phantom}


\section{序}
\label{more-groupoids-section-introduction}

\noindent
本章では亜群スキームに関する進んだ話題を扱う。
結果は亜群スキームの言葉で述べるが、読者は $2$-Cartesian 図式
\begin{equation}
\label{more-groupoids-equation-quotient-stack}
\vcenter{
\xymatrix{
R \ar[r] \ar[d] & U \ar[d] \\
U \ar[r] & [U/R]
}
}
\end{equation}
を念頭に置くとよい。ここで $[U/R]$ は商スタックである。『代数空間における亜群』
注意 \ref{spaces-groupoids-remark-fundamental-square} を参照せよ。
多くの結果はこの図式を考察することから動機づけられる。
例えば Keel と Mori の美しい論文 \cite{K-M} を参照されたい。





\section{記法}
\label{more-groupoids-section-notation}

\noindent
引き続き『亜群スキーム』節 \ref{groupoids-section-notation} で導入した規約と
記法に従う。






\section{有用な図式}
\label{more-groupoids-section-diagrams}

\noindent
本章で容易に参照できるよう、『亜群スキーム』補題
\ref{groupoids-lemma-diagram} および \ref{groupoids-lemma-diagram-pull}
の結果を簡潔に再述する。
$S$ をスキームとし、$(U, R, s, t, c)$ を $S$ 上の亜群スキームとする。
可換図式
\begin{equation}
\label{more-groupoids-equation-diagram}
\vcenter{
\xymatrix{
& U & \\
R \ar[d]_s \ar[ru]^t &
R \times_{s, U, t} R
\ar[l]^-{\text{pr}_0} \ar[d]^{\text{pr}_1} \ar[r]_-c &
R \ar[d]^s \ar[lu]_t \\
U & R \ar[l]_t \ar[r]^s & U
}
}
\end{equation}
において、下の二つの正方形はファイバー積図式である。
さらに、上の三角形（実際には正方形）も Cartesian である。

\medskip\noindent
図式
\begin{equation}
\label{more-groupoids-equation-pull}
\vcenter{
\xymatrix{
R \times_{t, U, t} R
\ar@<1ex>[r]^-{\text{pr}_1} \ar@<-1ex>[r]_-{\text{pr}_0}
\ar[d]_{\text{pr}_0 \times c \circ (i, 1)} &
R \ar[r]^t \ar[d]^{\text{id}_R} &
U \ar[d]^{\text{id}_U} \\
R \times_{s, U, t} R
\ar@<1ex>[r]^-c \ar@<-1ex>[r]_-{\text{pr}_0} \ar[d]_{\text{pr}_1} &
R \ar[r]^t \ar[d]^s &
U \\
R \ar@<1ex>[r]^s \ar@<-1ex>[r]_t &
U
}
}
\end{equation}
は可換である。上の二つの行は、表示した垂直射により同型である。
左下の二つの正方形は Cartesian である。





\section{微分層}
\label{more-groupoids-section-differentials}

\noindent
次の補題は『亜群スキーム』補題
\ref{groupoids-lemma-group-scheme-module-differentials} に対応する結果である。

\begin{lemma}
\label{more-groupoids-lemma-sheaf-differentials}
$S$ をスキームとし、$(U, R, s, t, c)$ を $S$ 上の亜群スキームとする。
$t$ により $U$ 上のスキームとみなした $R$ の微分層は、はめ込み
$e : U \to R$ の余法層を $t$ で引き戻したものの商である。すなわち標準的全射
$t^*\mathcal{C}_{U/R} \to \Omega_{R/U}$ が存在する。$s$ が平坦ならば、
この写像は同型である。
\end{lemma}

\begin{proof}
$e : U \to R$ は射 $s$ の切断なので、はめ込みである。
『スキーム』補題 \ref{schemes-lemma-section-immersion} を参照せよ。
次の図式を考える。
$$
\xymatrix{
R \ar[r]_-{(1, i)} \ar[d]_t &
R \times_{s, U, t} R \ar[d]^c \ar[rr]_{(\text{pr}_0, i \circ \text{pr}_1)} & &
R \times_{t, U, t} R \\
U \ar[r]^e &
R
}
$$
左の正方形は Cartesian である。実際、$a \circ b = e$ ならば
$b = i(a)$ である。水平射の合成は $t : R \to U$ の対角射であり、
右上の水平射は同型である。したがって $\Omega_{R/U}$ はこの合成の余法層なので、
$(1, i)$ の余法層と同型である。『射』補題
\ref{morphisms-lemma-conormal-functorial-flat} により全射
$t^*\mathcal{C}_{U/R} \to \Omega_{R/U}$ を得る。また $c$ が平坦ならば、
これは同型である。図式 (\ref{more-groupoids-equation-pull}) の性質により $c$ は $s$ の
基底変換だから、$s$ が平坦ならば $c$ も平坦である。
『射』補題 \ref{morphisms-lemma-base-change-flat} を参照せよ。
\end{proof}






\section{局所構造}
\label{more-groupoids-section-local}

\noindent
$S$ をスキームとし、$(U, R, s, t, c, e, i)$ を $S$ 上の亜群スキームとする。
$u \in U$ を点とする。本節では局所環
$$
A = \mathcal{O}_{U, u}
\quad\text{and}\quad
B = \mathcal{O}_{R, e(u)}
$$
上に得られる構造を説明する。射 $s, t, c, e, i$ が誘導する局所環準同型を、
対応する同じ文字で表すことにする。特に、局所環の可換図式
$$
\xymatrix{
A \ar[rd]_t \ar[rrd]^1 \\
& B \ar[r]^e & A \\
A \ar[ru]^s \ar[rru]_1
}
$$
を得る。したがって $I \subset B$ を $e : B \to A$ の核とすれば、
$B = s(A) \oplus I = t(A) \oplus I$ である。また
$$
C = \mathcal{O}_{R \times_{s, U, t} R, (e(u), e(u))}
$$
とおく。このとき
$$
C = (B \otimes_{s, A, t} B)_{\mathfrak m_B \otimes B + B \otimes \mathfrak m_B}
$$
である。$J \subset C$ を $I \otimes B + B \otimes I$ が生成する $C$ のイデアルと
する。すると $J$ は局所環準同型
$$
(e, e) : C \longrightarrow A
$$
の核でもある。合成則 $c : R \times_{s, U, t} R \to R$ は、$I$ を $J$ に送る環準同型
$$
c : B \longrightarrow C
$$
に対応する。

\begin{lemma}
\label{more-groupoids-lemma-first-order-structure-c}
$c$ が誘導する写像 $I/I^2 \to J/J^2$ は合成
$$
I/I^2 \xrightarrow{(1, 1)} I/I^2 \oplus I/I^2 \to J/J^2
$$
である。ここで第二の矢印は等式 $J = (I \otimes B + B \otimes I)C$ から得られる。
写像 $i : B \to B$ は写像 $-1 : I/I^2 \to I/I^2$ を誘導する。
\end{lemma}

\begin{proof}
$C$ から別の局所環への局所準同型を記述するには、$b_1 \otimes b_2$ の形の元に
何が起こるかを述べれば十分である。これを踏まえると、二つの標準的写像
$$
e_2 : C \to B,\ b_1 \otimes b_2 \mapsto b_1s(e(b_2)),\quad
e_1 : C \to B,\ b_1 \otimes b_2 \mapsto t(e(b_1))b_2
$$
がある。これらは、それぞれ $r \mapsto (r, e(s(r)))$ および
$r \mapsto (e(t(r)), r)$ で与えられる埋め込み
$R \to R \times_{s, U, t} R$ に対応する。
これらの写像が定める写像 $J/J^2 \to I/I^2$ は、合わせて補題の写像
$I/I^2 \oplus I/I^2 \to J/J^2$ の逆写像を与える。したがって示すべきことは、
$e_1 \circ c : B \to B$ と $e_2 \circ c : B \to B$ が恒等写像であることだけである。
これは二つの合成 $R \to R \times_{s, U, t} R \to R$ がとも恒等写像であることから
従う。

\medskip\noindent
$i$ に関する主張は、$c$ に関する主張と等式 $c \circ (1, i) = e \circ t$ から従う。
詳細の一部は省略する。
\end{proof}






\section{亜群の性質}
\label{more-groupoids-section-technical-lemma}

\noindent
$(U, R, s, t, c)$ を亜群スキームとする。
本節の結果の背後にある考えは、$s: R \to U$ が射 $U \to [U/R]$ の
基底変換であるということである（図式 (\ref{more-groupoids-equation-quotient-stack}) を参照）。
したがって、$s : R \to U$ の局所的性質は射 $U \to [U/R]$ の局所的性質を
反映するはずである。しかし $[U/R]$ は常に代数スタックであるとは限らず、
したがって $U \to [U/R]$ の幾何学的または代数的性質を語れないため、
この考えをそのまま用いることはできない。それでも、商スタックに一切
言及せずに、この考えの一部を実現できることが分かる。

\medskip\noindent
そのような結果の最初の例を示す。大まかに言えば、補題で得られる開集合
$W \subset U'$ は、射 $U' \to [U/R]$ が性質 $\mathcal{P}$ をもつ点の軌跡である。

\begin{lemma}
\label{more-groupoids-lemma-local-source}
$S$ をスキームとし、$(U, R, s, t, c, e, i)$ を $S$ 上の亜群とする。
$g : U' \to U$ をスキームの射とし、合成
$$
\xymatrix{
h : U' \times_{g, U, t} R \ar[r]_-{\text{pr}_1} & R \ar[r]_s & U.
}
$$
を $h$ と書く。$\mathcal{P}, \mathcal{Q}, \mathcal{R}$ をスキームの射の性質とし、
次を仮定する。
\begin{enumerate}
\item $\mathcal{R} \Rightarrow \mathcal{Q}$ である。
\item $\mathcal{Q}$ は基底変換および合成で保たれる。
\item $\mathcal{Q}$ をもつ任意の射 $f : X \to Y$ に対し、
$f|_{W(\mathcal{P}, f)}$ が $\mathcal{P}$ をもつような最大の開集合
$W(\mathcal{P}, f) \subset X$ が存在する。
\item $\mathcal{Q}$ をもつ任意の射 $f : X \to Y$ と、$\mathcal{R}$ をもつ任意の射
$Y' \to Y$ に対して $Y' \times_Y W(\mathcal{P}, f) = W(\mathcal{P}, f')$ が成り立つ。
ここで $f' : X_{Y'} \to Y'$ は $f$ の基底変換である。
\end{enumerate}
$s, t$ が $\mathcal{R}$ をもち、$g$ が $\mathcal{Q}$ をもつならば、
$W \times_{g, U, t} R = W(\mathcal{P}, h)$ を満たす開部分スキーム
$W \subset U'$ が存在する。
\end{lemma}

\begin{proof}
次の図式は可換であり、二つの正方形はいずれも Cartesian である。
$$
\xymatrix{
U' \times_{g, U, t} R \times_{t, U, t} R
\ar[rr]_-{\text{pr}_{12}}
\ar@<1ex>[d]^-{\text{pr}_{02}} \ar@<-1ex>[d]_-{\text{pr}_{01}} & &
R \times_{t, U, t} R
\ar@<1ex>[d]^-{\text{pr}_1} \ar@<-1ex>[d]_-{\text{pr}_0}
\\
U' \times_{g, U, t} R \ar[rr]^{\text{pr}_1} & & R
}
$$
ここでは二つの写像 $t \circ \text{pr}_i : R \times_{t, U, t} R \to U$ が
等しいことを用いた。これと図式 (\ref{more-groupoids-equation-pull}) の性質を合わせると、
両方の正方形が Cartesian である可換図式
$$
\xymatrix{
U' \times_{g, U, t} R \times_{t, U, t} R
\ar[rr]_-{c \circ (i, 1)}
\ar@<1ex>[d]^-{\text{pr}_{02}} \ar@<-1ex>[d]_-{\text{pr}_{01}} & &
R
\ar@<1ex>[d]^-{s} \ar@<-1ex>[d]_-{t}
\\
U' \times_{g, U, t} R \ar[rr]^h & & U
}
$$
を得る。

\medskip\noindent
$s, t$ が $\mathcal{R}$ をもち、$g$ が $\mathcal{Q}$ をもつと仮定する。
$h$ は $s$（$\mathcal{R}$ をもち、したがって $\mathcal{Q}$ をもつ）と、
$g$ の基底変換（$\mathcal{Q}$ をもつ）との合成なので、$\mathcal{Q}$ をもつ。
よって $W(\mathcal{P}, h) \subset U' \times_{g, U, t} R$ が存在する。
仮定により
$\text{pr}_{01}^{-1}(W(\mathcal{P}, h)) =
\text{pr}_{02}^{-1}(W(\mathcal{P}, h))$
である。実際、どちらも $c \circ (i, 1)$ が $\mathcal{P}$ をもつ最大の開集合である。
射影 $U' \times_{g, U, t} R \to U'$ は切断
$\sigma : U' \to U' \times_{g, U, t} R$, $u' \mapsto (u', e(g(u')))$
をもつことに注意する。また同型
$$
(U' \times_{g, U, t} R) \times_{U'} (U' \times_{g, U, t} R)
=
U' \times_{g, U, t} R \times_{t, U, t} R
$$
を通じて、左辺から $U' \times_{g, U, t} R$ への二つの射影は、右辺の射
$\text{pr}_{01}$ および $\text{pr}_{02}$ と一致する。したがって
$\text{pr}_{01}^{-1}(W(\mathcal{P}, h)) =
\text{pr}_{02}^{-1}(W(\mathcal{P}, h))$
より、$W(\mathcal{P}, h)$ は $U$ のある部分集合の逆像である。この部分集合は必然的に
開集合 $W = \sigma^{-1}(W(\mathcal{P}, h))$ である。
\end{proof}

\begin{remark}
\label{more-groupoids-remark-local-source-warning}
注意：補題 \ref{more-groupoids-lemma-local-source} は慎重に用いる必要がある。
例えば、この補題は $\mathcal{P}=$``flat'', $\mathcal{Q}=$``empty'',
$\mathcal{R}=$``平坦かつ局所有限表示'' の場合に適用できる。
しかしスキームの射 $f : X \to Y$ に対し、$f|_W$ が平坦となる最大の開集合
$W \subset X$ は、$f$ が平坦である点の集合では {\it ない}！
\end{remark}

\begin{remark}
\label{more-groupoids-remark-local-source-apply}
注意 \ref{more-groupoids-remark-local-source-warning} にもかかわらず、補題
\ref{more-groupoids-lemma-local-source} を過度な曖昧さなしに用いられる場合がある。
以下にその一覧を挙げる。各場合について、仮定 (1) と (2) の確認は省略し、
(3) と (4) を導く参照先を示す。
\begin{enumerate}
\item $\mathcal{Q} = \mathcal{R} =$``locally of finite type'', かつ
$\mathcal{P} =$``relative dimension $\leq d$'' の場合。
『射』定義 \ref{morphisms-definition-relative-dimension-d} および補題
\ref{morphisms-lemma-openness-bounded-dimension-fibres},
\ref{morphisms-lemma-dimension-fibre-after-base-change} を参照せよ。
\item $\mathcal{Q} = \mathcal{R} =$``locally of finite type'', かつ
$\mathcal{P} =$``locally quasi-finite'' の場合。
これは前項で $d = 0$ とした場合である。『射』補題
\ref{morphisms-lemma-locally-quasi-finite-rel-dimension-0} を参照せよ。
\item $\mathcal{Q} = \mathcal{R} =$``locally of finite type'', かつ
$\mathcal{P} =$``unramified'' の場合。
『射』補題 \ref{morphisms-lemma-unramified-characterize} および
\ref{morphisms-lemma-set-points-where-fibres-unramified} を参照せよ。
\end{enumerate}
上に挙げた場合で興味深いのは、射 $h$ が性質 $\mathcal{P}$ をもつ軌跡について
結論を得るために、$s, t$ が平坦であると仮定する必要がないことである。
一覧を続ける。
\begin{enumerate}
\item[(4)] $\mathcal{Q} =$``局所有限表示'',
$\mathcal{R} =$``平坦かつ局所有限表示'', かつ
$\mathcal{P} =$``flat'' の場合。『射の続論』定理
\ref{more-morphisms-theorem-openness-flatness} および補題
\ref{more-morphisms-lemma-flat-locus-base-change} を参照せよ。
\item[(5)] $\mathcal{Q} =$``局所有限表示'',
$\mathcal{R} =$``平坦かつ局所有限表示'', かつ
$\mathcal{P}=$``Cohen-Macaulay'' の場合。『射の続論』定義
\ref{more-morphisms-definition-CM} および補題
\ref{more-morphisms-lemma-base-change-CM},
\ref{more-morphisms-lemma-flat-finite-presentation-CM-open} を参照せよ。
\item[(6)] $\mathcal{Q} =$``局所有限表示'',
$\mathcal{R} =$``平坦かつ局所有限表示'', かつ
$\mathcal{P}=$``syntomic'' の場合は、『射』補題
\ref{morphisms-lemma-set-points-where-fibres-lci} を用いる
（この軌跡は自動的に開である）。
\item[(7)] $\mathcal{Q} =$``局所有限表示'',
$\mathcal{R} =$``平坦かつ局所有限表示'', かつ
$\mathcal{P}=$``smooth'' の場合。『射』補題
\ref{morphisms-lemma-set-points-where-fibres-smooth} を参照せよ
（この軌跡は自動的に開である）。
\item[(8)] $\mathcal{Q} =$``局所有限表示'',
$\mathcal{R} =$``平坦かつ局所有限表示'', かつ
$\mathcal{P}=$``\'etale'' の場合。『射』補題
\ref{morphisms-lemma-set-points-where-fibres-etale} を参照せよ
（この軌跡は自動的に開である）。
\end{enumerate}
\end{remark}

\noindent
第二の結果を示す。$R$-不変な開集合 $W \subset U$ は、射 $U \to [U/R]$ が
性質 $\mathcal{P}$ をもつような $[U/R]$ の最大の開集合の逆像と考えるべきである。

\begin{lemma}
\label{more-groupoids-lemma-property-invariant}
$S$ をスキームとし、$(U, R, s, t, c)$ を $S$ 上の亜群とする。
$\tau \in \{Zariski, \linebreak[0] fppf,
\linebreak[0] \etale, \linebreak[0]
smooth, \linebreak[0] syntomic\}$\footnote{$fpqc$ が含まれていないのは
誤植ではない。} とする。$\mathcal{P}$ を、終域上 $\tau$-局所的なスキームの射の
性質とする（『降下』定義 \ref{descent-definition-property-morphisms-local}）。
$\{s : R \to U\}$ と $\{t : R \to U\}$ が $\tau$-位相の被覆であると仮定する。
$s|_{s^{-1}(W)} : s^{-1}(W) \to W$ が性質 $\mathcal{P}$ をもつような最大の
開部分スキームを $W \subset U$ とする。このとき $W$ は $R$-不変である。
『亜群スキーム』定義 \ref{groupoids-definition-invariant-open} を参照せよ。
\end{lemma}

\begin{proof}
開集合 $W \subset U$ の存在と性質は、『降下』補題
\ref{descent-lemma-largest-open-of-the-base} に記述されている。
図式 (\ref{more-groupoids-equation-diagram}) において、射
$\text{pr}_1 : R \times_{s, U, t} R \to R$ が性質 $\mathcal{P}$ をもつような
最大の開部分スキームを $W_1 \subset R$ とする。前述の『降下』補題
\ref{descent-lemma-largest-open-of-the-base} と、$\{s : R \to U\}$ および
$\{t : R \to U\}$ が $\tau$-位相の被覆であるという仮定から、望みどおり
$t^{-1}(W) = W_1 = s^{-1}(W)$ を得る。
\end{proof}

\begin{lemma}
\label{more-groupoids-lemma-property-G-invariant}
$S$ をスキームとし、$(U, R, s, t, c)$ を $S$ 上の亜群とする。
$G \to U$ をその安定化群スキームとする。
$\tau \in \{fppf, \linebreak[0] \etale, \linebreak[0]
smooth, \linebreak[0] syntomic\}$ とし、$\mathcal{P}$ を終域上 $\tau$-局所的な
射の性質とする。$\{s : R \to U\}$ と $\{t : R \to U\}$ が $\tau$-位相の
被覆であると仮定する。$G_W \to W$ が性質 $\mathcal{P}$ をもつような最大の
開部分スキームを $W \subset U$ とする。このとき $W$ は $R$-不変である
（『亜群スキーム』定義 \ref{groupoids-definition-invariant-open} を参照）。
\end{lemma}

\begin{proof}
開集合 $W \subset U$ の存在と性質は、『降下』補題
\ref{descent-lemma-largest-open-of-the-base} に記述されている。射
$$
G \times_{U, t} R \longrightarrow R \times_{s, U} G, \quad
(g, r) \longmapsto (r, r^{-1} \circ g \circ r)
$$
は $R$ 上の同型である（ここで $\circ$ は亜群における合成を表す）。
したがって、前述の『降下』補題
\ref{descent-lemma-largest-open-of-the-base} で示された $W$ の性質により、
$s^{-1}(W) = t^{-1}(W)$ である。
\end{proof}



\section{ファイバーの比較}
\label{more-groupoids-section-fibres}

\noindent
$(U, R, s, t, c, e, i)$ を $S$ 上の亜群スキームとする。
図式 (\ref{more-groupoids-equation-diagram}) により、亜群における写像 $s : R \to U$ の
ファイバーを比較できる。点 $u \in U$ に対し、$s : R \to U$ の $u$ 上の
スキーム論的ファイバーを $F_u = s^{-1}(u)$ と書く。例えばこの図式から、
$u, u' \in U$ が $s(r) = u$, $t(r) = u'$ を満たす点ならば
$(F_u)_{\kappa(r)} \cong (F_{u'})_{\kappa(r)}$ であることが従う。
これは、以下のより一般的で精密な補題 \ref{more-groupoids-lemma-two-fibres} の特別な場合である。
実際、$r' = i(r)$ とすればよい。

\medskip\noindent
スキーム $X$ と点 $x \in X$ の組 $(X, x)$ を、$x$ における $X$ の {\it 芽} と
呼ぶことがある。{\it 芽の射} $f : (X, x) \to (S, s)$ とは、$x$ のある開近傍上で
定義され、$f(x) = s$ を満たす射 $f : U \to S$ のことである。このような二つの射
$f$, $f'$ が同じ芽の射を与えるとは、$x$ のある開近傍で $f$ と $f'$ が一致する
こととする。$\tau \in \{Zariski, \etale, smooth, syntomic, fppf\}$ とする。
ここで一時的に次の概念を導入する。二つの芽の射
$f : (X, x) \to (S, s)$ と $f' : (X', x') \to (S', s')$ が
{\it 基底上 $\tau$-位相で局所的に同型} であるとは、点付きスキーム
$(S'', s'')$ と芽の射 $g : (S'', s'') \to (S, s)$,
$g' : (S'', s'') \to (S', s')$ が存在し、次を満たすこととする。
\begin{enumerate}
\item $g$ と $g'$ は $s''$ において開はめ込み（それぞれ\ \'etale、smooth、
syntomic、または平坦かつ局所有限表示）である。
\item $(s'', x)$ および $(s'', x')$ の上にある点 $\tilde x$, $\tilde x'$ を適当に
選ぶと、芽 $(S'', s'')$ 上の芽の同型
$$
(S'' \times_{g, S, f} X, \tilde x)
\cong
(S'' \times_{g', S', f'} X', \tilde  x')
$$
が存在する。
\end{enumerate}
最後に、上と同様の $S'', s'', g, g'$ が存在し、(1) の代わりに $g$ と $g'$ が
$s''$ において平坦であるという条件を課せるとき、芽の射
$f : (X, x) \to (S, s)$ と $f' : (X', x') \to (S', s')$ は
{\it 基底上平坦局所的に同型} であるという。この条件は上のどの $\tau$ 条件よりも
はるかに弱い。実際、平坦射は開射であるとは限らない。

\begin{lemma}
\label{more-groupoids-lemma-two-fibres}
$S$ をスキームとし、$(U, R, s, t, c)$ を $S$ 上の亜群とする。
$r, r' \in R$ が $U$ において $t(r) = t(r')$ を満たすとし、
$u = s(r)$, $u' = s(r')$ とおく。スキーム論的ファイバーを
$F_u = s^{-1}(u)$, $F_{u'} = s^{-1}(u')$ と書く。
\begin{enumerate}
\item 共通の体拡大 $\kappa(u) \subset k$, $\kappa(u') \subset k$ と同型
$(F_u)_k \cong (F_{u'})_k$ が存在する。
\item (1) の同型は、$r$ の上にある点を $r'$ の上にある点へ写すように選べる。
\item 射 $s$, $t$ が平坦ならば、芽の射
$s : (R, r) \to (U, u)$ と $s : (R, r') \to (U, u')$ は基底上平坦局所的に同型である。
\item 射 $s$, $t$ が \'etale（それぞれ\ smooth、syntomic、または平坦かつ
局所有限表示）ならば、芽の射 $s : (R, r) \to (U, u)$ と
$s : (R, r') \to (U, u')$ は、\'etale（それぞれ\ smooth、syntomic、または
fppf）位相で基底上局所的に同型である。
\end{enumerate}
\end{lemma}

\begin{proof}
図式 (\ref{more-groupoids-equation-diagram}) の存在と性質を繰り返し用いる。
この図式の性質と『スキーム』補題 \ref{schemes-lemma-points-fibre-product} により、
$\text{pr}_0(\xi) = r$, $c(\xi) = r'$ を満たす
$R \times_{s, U, t} R$ の点 $\xi$ が存在する。
$\tilde r = \text{pr}_1(\xi) \in R$ とおく。

\medskip\noindent
(1) の証明。$k = \kappa(\tilde r)$ とおく。$t(\tilde r) = u$ かつ
$s(\tilde r) = u'$ なので、$k$ は $\kappa(u)$ と $\kappa(u')$ の共通拡大である。
さらに $(F_u)_k$ と $(F_{u'})_k$ はとも、$\tilde r$ 上の
$\text{pr}_1 : R \times_{s, U, t} R \to R$ のファイバーと同型である。
よって (1) が示された。

\medskip\noindent
点 $\xi$ はそれぞれ\ $r$, $r'$ に写るので、(2) が従う。

\medskip\noindent
(3) は、$\tilde u$ と $\tilde u'$ に点 $\xi$ を用いれば、上の議論と定義から明らかである。

\medskip\noindent
$s$ と $t$ が平坦かつ有限表示ならば、これらは開射である
（『射』補題 \ref{morphisms-lemma-fppf-open}）。したがって $\tilde r$ のある
アフィン開近傍 $V''$ の像は、$u$ の開近傍 $V$ と、それぞれ\ $u'$ の開近傍 $V'$ を
被覆する。これらを用いれば、「基底上 $\tau$-位相で局所的に同型」の定義における
性質 (1) と (2) を示せる。
\end{proof}







\section{Cohen--Macaulay 表示}
\label{more-groupoids-section-CM}

\noindent
$s, t$ が平坦かつ局所有限表示である任意の亜群 $(U, R, s, t, c)$ に対し、
$s'$ と $t'$ が Cohen--Macaulay 射（かつ局所有限表示）となる「同値な」亜群
$(U', R', s', t', c')$ が存在する。Cohen--Macaulay 射について詳しくは
『射の続論』節 \ref{more-morphisms-section-CM} を参照せよ。ここで「同値」とは、
商スタック $[U/R]$ と $[U'/R']$ がスタックとして同値であることと解釈できる。
『代数空間における亜群』節 \ref{spaces-groupoids-section-stacks} および
節 \ref{spaces-groupoids-section-quotient-stack-restrict} を参照せよ。

\begin{lemma}
\label{more-groupoids-lemma-make-CM}
$S$ をスキームとし、$(U, R, s, t, c)$ を $S$ 上の亜群とする。
$s$ と $t$ が平坦かつ局所有限表示であると仮定する。このとき、次を満たす開集合
$U' \subset U$ が存在する。
\begin{enumerate}
\item $t^{-1}(U') \subset R$ は、射 $s$ が Cohen--Macaulay となる $R$ の
最大の開部分スキームである。
\item $s^{-1}(U') \subset R$ は、射 $t$ が Cohen--Macaulay となる $R$ の
最大の開部分スキームである。
\item 射 $t|_{s^{-1}(U')} : s^{-1}(U') \to U$ は全射である。
\item 射 $s|_{t^{-1}(U')} : t^{-1}(U') \to U$ は全射である。
\item $R$ を $U'$ に制限した $R' = s^{-1}(U') \cap t^{-1}(U')$ は亜群
$(U', R', s', t', c')$ を定め、射 $s'$ と $t'$ は Cohen--Macaulay かつ
局所有限表示である。
\end{enumerate}
\end{lemma}

\begin{proof}
補題 \ref{more-groupoids-lemma-local-source} を、$g = \text{id}$ および
$\mathcal{Q} =$``局所有限表示'',
$\mathcal{R} =$``平坦かつ局所有限表示'',
$\mathcal{P}=$``Cohen-Macaulay'' として適用する。注意
\ref{more-groupoids-remark-local-source-apply} を参照せよ。これにより、$t^{-1}(U') \subset R$ が
$R$ の中で射 $s$ の Cohen--Macaulay 軌跡である最大の開部分スキームとなるような
開集合 $U' \subset U$ を得る。これで (1) が示された。
$i : R \to R$ を亜群の逆元写像とする。$i$ は同型であり、
$s \circ i = t$, $t \circ i = s$ なので、$s^{-1}(U')$ も射 $t$ が
Cohen--Macaulay となる $R$ の最大の開集合である。これで (2) が示された。
『射の続論』補題 \ref{more-morphisms-lemma-flat-finite-presentation-CM-open} により、
開部分集合 $t^{-1}(U')$ は $s : R \to U$ の各ファイバーで稠密である。
これで (3) が示された。(4) も同じ議論による。
(5) は、(1), (2) と『亜群スキーム』節
\ref{groupoids-section-restrict-groupoid} における制限の議論から形式的に従う。
\end{proof}








\section{亜群の制限}
\label{more-groupoids-section-restricting-groupoids}

\noindent
本節では、制限に継承される亜群の性質に関する補題をまとめる。
これらの補題の大半は、制限を定義する図式
\begin{equation}
\label{more-groupoids-equation-restriction}
\vcenter{
\xymatrix{
R' \ar[d] \ar[r] \ar@/_3pc/[dd]_{t'} \ar@/^1pc/[rr]^{s'}&
R \times_{s, U} U' \ar[r] \ar[d] &
U' \ar[d]^g \\
U' \times_{U, t} R \ar[d] \ar[r] &
R \ar[r]^s \ar[d]_t &
U \\
U' \ar[r]^g &
U
}
}
\end{equation}
を考察すれば証明できる。『亜群スキーム』補題
\ref{groupoids-lemma-restrict-groupoid} を参照せよ。

\begin{lemma}
\label{more-groupoids-lemma-restrict-preserves-type}
$S$ をスキームとし、$(U, R, s, t, c)$ を $S$ 上の亜群スキームとする。
$g : U' \to U$ をスキームの射とし、$(U', R', s', t', c')$ を $g$ による
$(U, R, s, t, c)$ の制限とする。
\begin{enumerate}
\item $s, t$ と $g$ が局所有限型ならば、$s', t'$ は局所有限型である。
\item $s, t$ と $g$ が局所有限表示ならば、$s', t'$ は局所有限表示である。
\item $s, t$ と $g$ が平坦ならば、$s', t'$ は平坦である。
\item ここにさらに追加する。
\end{enumerate}
\end{lemma}

\begin{proof}
局所有限型という性質は合成および任意の基底変換で保たれる。
『射』補題 \ref{morphisms-lemma-composition-finite-type} および
\ref{morphisms-lemma-base-change-finite-type} を参照せよ。
したがって (1) は図式 (\ref{more-groupoids-equation-restriction}) から明らかである。
他の場合については、『射』補題
\ref{morphisms-lemma-composition-finite-presentation},
\ref{morphisms-lemma-base-change-finite-presentation},
\ref{morphisms-lemma-composition-flat},
\ref{morphisms-lemma-base-change-flat} を参照せよ。
\end{proof}

\noindent
次の補題を用いれば、前の補題の結果をより一様に証明することもできた。

\begin{lemma}
\label{more-groupoids-lemma-restrict-property}
$S$ をスキームとし、$(U, R, s, t, c)$ を $S$ 上の亜群スキームとする。
$g : U' \to U$ をスキームの射とし、$(U', R', s', t', c')$ を $g$ による
$(U, R, s, t, c)$ の制限とする。また
$h = s \circ \text{pr}_1 : U' \times_{g, U, t} R \to U$ とおく。
$\mathcal{P}$ がスキームの射の性質であり、
\begin{enumerate}
\item $h$ は性質 $\mathcal{P}$ をもち、
\item $\mathcal{P}$ は基底変換で保たれる
\end{enumerate}
ならば、$s', t'$ は性質 $\mathcal{P}$ をもつ。
\end{lemma}

\begin{proof}
図式 (\ref{more-groupoids-equation-restriction}) により $s'$ は $h$ の基底変換であり、
$t'$ はスキームの射として $s'$ と同型なので明らかである。
\end{proof}

\begin{lemma}
\label{more-groupoids-lemma-double-restrict}
$S$ をスキームとし、$(U, R, s, t, c)$ を $S$ 上の亜群スキームとする。
$g : U' \to U$, $g' : U'' \to U'$ をスキームの射とし、$g'' = g \circ g'$ とおく。
$(U', R', s', t', c')$ を $R$ の $U'$ への制限とする。また
$h = s \circ \text{pr}_1 : U' \times_{g, U, t} R \to U$,
$h' = s' \circ \text{pr}_1 : U'' \times_{g', U', t} R \to U'$,
$h'' = s \circ \text{pr}_1 : U'' \times_{g'', U, t} R \to U$ とおく。
次の図式は可換であり、両方の正方形は Cartesian である。
$$
\xymatrix{
U'' \times_{g', U', t} R' \ar[d]^{h'} &
(U' \times_{g, U, t} R) \times_U (U'' \times_{g'', U, t} R)
\ar[l] \ar[r] \ar[d] &
U'' \times_{g'', U, t} R \ar[d]_{h''} \\
U' &
U' \times_{g, U, t} R \ar[l]_{\text{pr}_0} \ar[r]^h &
U
}
$$
ここで左上の水平射は、証明で説明する記法を用いて次の規則で与えられる。
$$
\begin{matrix}
(U' \times_{g, U, t} R) \times_U (U'' \times_{g'', U, t} R) &
\longrightarrow &
U'' \times_{g', U', t} R' \\
((u', r_0), (u'', r_1)) &
\longmapsto &
(u'', (c(r_1, i(r_0)), (g'(u''), u')))
\end{matrix}
$$
\end{lemma}

\begin{proof}
関手的な見方を利用し、この補題を亜群圏の制限における矢印についての主張へ
帰着して具体的に確かめる。補題の最後の式における記法
$((u', r_0), (u'', r_1))$ は、
$(U' \times_{g, U, t} R) \times_U (U'' \times_{g'', U, t} R)$ の
$T$-値点を表す。これは $u', u'', r_0, r_1$ がそれぞれ $U', U'', R, R$ の
$T$-値点であり、$g(u') = t(r_0)$,
$g(g'(u'')) = g''(u'') = t(r_1)$, $s(r_0) = s(r_1)$ が成り立つことを意味する。
厳密には $g \circ u' = t \circ r_0$ などと書くべきだが、そうすると記法が
さらに読みにくくなる。$r_1$ と $r_0$ を亜群圏の矢印と考えると、これは図式
$$
\xymatrix{
t(r_0) = g(u') &
s(r_0) = s(r_1) \ar[l]_{r_0} \ar[r]^-{r_1} &
t(r_1) = g(g'(u''))
}
$$
で表せる。特に、この図式は合成 $c(r_1, i(r_0))$ が意味をもつことを示す。
ここで
$$
R' = R \times_{(t, s), U \times_S U, g \times g} U' \times_S U'
$$
を思い出すと、$R'$ の $T$-値点は $(r, (u'_0, u'_1))$ の形であり、
$t(r) = g(u'_0)$ と $s(r) = g(u'_1)$ を満たす。したがって、上の
$((u', r_0), (u'', r_1))$ から $R'$ の $T$-値点
$(c(r_1, i(r_0)), (g'(u''), u'))$ を得る。実際、
$t(c(r_1, i(r_0))) = t(r_1) = g(g'(u''))$ および
$s(c(r_1, i(r_0))) = s(i(r_0)) = t(r_0) = g(u')$ が成り立つ。
この定義によって左の正方形が可換となることの確認は読者に委ねる。

\medskip\noindent
左の正方形が Cartesian であることを示す。$v' = s'(p')$ を満たす
$U'' \times_{g', U', t} R'$ と $U' \times_{g, U, t} R$ の $T$-値点
$(v'', p')$, $(v', p)$ が与えられたとする。このとき
$g'(v'') = t'(p')$ および $g(v') = t(p)$ でもある。上の議論により、
$p' = (r, (u_0', u_1'))$ と書けて、$t(r) = g(u'_0)$,
$s(r) = g(u'_1)$ が成り立つ。この記法のもとで
$v' = s'(p') = u_1'$, $g'(v'') = t'(p') = u_0'$ である。図式で表すと
$$
\xymatrix{
s(p) \ar[r]^-p &
g(v') = g(u'_1) \ar[r]^-r &
g(u'_0) = g(g'(v''))
}
$$
となる。示すべきことは、$v' = u'$, $p = r_0$, $v'' = u''$ および
$p' = (c(r_1, i(r_0)), (g'(u''), u'))$ を満たす、上の形の一意な $T$-値点
$((u', r_0), (u'', r_1))$ が存在することである。上の二つの図式を比較すれば、
次を選ぶほかないことは明らかである。
$$
((u', r_0), (u'', r_1)) = ((v', p), (v'', c(r, p))
$$
詳細の一部は省略する。
\end{proof}

\begin{lemma}
\label{more-groupoids-lemma-double-restrict-property}
$S$ をスキームとし、$(U, R, s, t, c)$ を $S$ 上の亜群スキームとする。
$g : U' \to U$, $g' : U'' \to U'$ をスキームの射とし、$g'' = g \circ g'$ とおく。
$(U', R', s', t', c')$ を $R$ の $U'$ への制限とする。また
$h = s \circ \text{pr}_1 : U' \times_{g, U, t} R \to U$,
$h' = s' \circ \text{pr}_1 : U'' \times_{g', U', t} R \to U'$,
$h'' = s \circ \text{pr}_1 : U'' \times_{g'', U, t} R \to U$ とおく。
$\tau \in \{Zariski, \linebreak[0] \etale, \linebreak[0]
smooth, \linebreak[0] syntomic, \linebreak[0] fppf, \linebreak[0] fpqc\}$ とする。
$\mathcal{P}$ を、基底変換で保たれ、終域上 $\tau$-位相に関して局所的な
スキームの射の性質とする。もし
\begin{enumerate}
\item $h(U' \times_U R)$ は $U$ で開であり、
\item $\{h : U' \times_U R \to h(U' \times_U R)\}$ は $\tau$-被覆であり、
\item $h'$ は性質 $\mathcal{P}$ をもつ
\end{enumerate}
ならば、$h''$ は性質 $\mathcal{P}$ をもつ。逆に、もし
\begin{enumerate}
\item[(a)] $\{t : R \to U\}$ は $\tau$-被覆であり、
\item[(d)] $h''$ は性質 $\mathcal{P}$ をもつ
\end{enumerate}
ならば、$h'$ は性質 $\mathcal{P}$ をもつ。
\end{lemma}

\begin{proof}
補題 \ref{more-groupoids-lemma-double-restrict} の図式の性質から形式的に従う。
最初の場合、$g'' = g \circ g'$ なので、射 $h''$ の像は $h$ の像に含まれる。
したがって図式の右下隅の $U$ を $h(U' \times_U R)$ で置き換えられる。
これが補題の条件 (1), (2) の意味を説明する。第二の場合、
$\{\text{pr}_0 : U' \times_{g, U, t} R \to U'\}$ は、$\tau$ の基底変換と
条件 (a) により $\tau$-被覆であることに注意する。
\end{proof}






\section{体上の亜群の性質}
\label{more-groupoids-section-properties-groupoids-on-fields}

\noindent
「体上の亜群」とは、$U$ がある体のスペクトルである亜群スキーム
$(U, R, s, t, c)$ をいう。$(U, R, s, t, c)$ が体上で定義されているという
意味では {\bf なく}、より正確には、射 $s, t : R \to U$ が等しいという意味では
{\bf ない}。任意の体 $k$、抽象群 $G$、群準同型
$\varphi : G \to \text{Aut}(k)$ に対し、次のように定めることで
$\mathbf{Z}$ 上の亜群スキーム $(U, R, s, t, c)$ を得る。
\begin{align*}
U & = \Spec(k) \\
R & = \coprod\nolimits_{g \in G} \Spec(k) \\
s & = \coprod\nolimits_{g \in G} \Spec(\text{id}_k) \\
t & = \coprod\nolimits_{g \in G} \Spec(\varphi(g)) \\
c & = \text{composition in }G
\end{align*}
この例はなお $\Spec(k^G)$ 上の亜群スキームである。したがって $G$ が有限ならば、
$U = \Spec(k)$ は $\Spec(k^G)$ 上有限である。ある意味で本節の目標は、$s, t$ に
適切な有限性条件を課すと、体上の任意の亜群が有限指数部分体
$k' \subset k$ 上で定義されることを示すことである。

\medskip\noindent
$k$ を体とし、$(G, m)$ を $k$ 上の群スキーム、$p : G \to \Spec(k)$ をその
構造射とする。このとき $(\Spec(k), G, p, p, m)$ は体上の亜群の例である
（この場合はもちろん構造全体が体上で定義されている）。したがって本節は
『亜群スキーム』節 \ref{groupoids-section-properties-group-schemes-field} の
類似とみなせる。

\begin{lemma}
\label{more-groupoids-lemma-groupoid-on-field-open-multiplication}
$S$ をスキームとし、$(U, R, s, t, c)$ を $S$ 上の亜群スキームとする。
$U$ が体のスペクトルならば、合成射
$c : R \times_{s, U, t} R \to R$ は開である。
\end{lemma}

\begin{proof}
図式 (\ref{more-groupoids-equation-pull}) により、合成は射影
$\text{pr}_1 : R \times_{t, U, t} R \to R$ と同型である。この射影は
『射』補題 \ref{morphisms-lemma-scheme-over-field-universally-open} により開である。
\end{proof}

\begin{lemma}
\label{more-groupoids-lemma-groupoid-on-field-separated}
$S$ をスキームとし、$(U, R, s, t, c)$ を $S$ 上の亜群スキームとする。
$U$ が体のスペクトルならば、$R$ は分離スキームである。
\end{lemma}

\begin{proof}
『亜群スキーム』補題 \ref{groupoids-lemma-group-scheme-over-field-separated} により、
安定化群スキーム $G \to U$ は分離である。『亜群スキーム』補題
\ref{groupoids-lemma-diagonal} により、射
$j = (t, s) : R \to U \times_S U$ は分離である。$U$ は体のスペクトルなので、
$U \times_S U$ はアフィンである（『スキーム』節
\ref{schemes-section-fibre-products} のファイバー積の構成による）。したがって
$R$ は分離スキームである。『スキーム』補題
\ref{schemes-lemma-separated-permanence} を参照せよ。
\end{proof}

\begin{lemma}
\label{more-groupoids-lemma-groupoid-on-field-homogeneous}
$S$ をスキームとし、$(U, R, s, t, c)$ を $S$ 上の亜群スキームとする。
$U = \Spec(k)$ で $k$ が体であると仮定する。任意の点 $r, r' \in R$ に対し、
体拡大 $k'/k$、点 $r_1, r_2 \in R \times_{s, \Spec(k)} \Spec(k')$、および図式
$$
\xymatrix{
R &
R \times_{s, \Spec(k)} \Spec(k')
\ar[l]_-{\text{pr}_0} \ar[r]^\varphi &
R \times_{s, \Spec(k)} \Spec(k')
\ar[r]^-{\text{pr}_0} &
R
}
$$
が存在し、$\varphi$ は $\Spec(k')$ 上のスキームの同型で、
$\varphi(r_1) = r_2$, $\text{pr}_0(r_1) = r$, $\text{pr}_0(r_2) = r'$ を満たす。
\end{lemma}

\begin{proof}
補題 \ref{more-groupoids-lemma-two-fibres} の (1), (2) の特別な場合である。
\end{proof}

\begin{lemma}
\label{more-groupoids-lemma-restrict-groupoid-on-field}
$S$ をスキームとし、$(U, R, s, t, c)$ を $S$ 上の亜群スキームとする。
$U = \Spec(k)$ で $k$ が体であると仮定する。$k'/k$ を体拡大、
$U' = \Spec(k')$ とし、$(U', R', s', t', c')$ を $U' \to U$ による
$(U, R, s, t, c)$ の制限とする。定義図式
$$
\xymatrix{
R' \ar[d] \ar[r] \ar@/_3pc/[dd]_{t'} \ar@/^1pc/[rr]^{s'} \ar@{..>}[rd] &
R \times_{s, U} U' \ar[r] \ar[d] &
U' \ar[d] \\
U' \times_{U, t} R \ar[d] \ar[r] &
R \ar[r]^s \ar[d]_t &
U \\
U' \ar[r] &
U
}
$$
のすべての射は全射、平坦、普遍的開である。さらに点線の射
$R' \to R$ はアフィンである。
\end{lemma}

\begin{proof}
射 $U' \to U$ は $\Spec(k') \to \Spec(k)$ に等しいので、アフィン、全射、平坦である。
射 $s, t : R \to U$ と $U' \to U$ は、『射』補題
\ref{morphisms-lemma-scheme-over-field-universally-open} により普遍的開である。
$R$ は空でなく、$U$ は体のスペクトルなので、射 $s, t : R \to U$ は全射かつ平坦である。
あとは『射』補題 \ref{morphisms-lemma-base-change-surjective},
\ref{morphisms-lemma-composition-surjective},
\ref{morphisms-lemma-composition-open},
\ref{morphisms-lemma-base-change-affine},
\ref{morphisms-lemma-composition-affine},
\ref{morphisms-lemma-base-change-flat},
\ref{morphisms-lemma-composition-flat} を用いれば結論を得る。
\end{proof}

\begin{lemma}
\label{more-groupoids-lemma-groupoid-on-field-explain-points}
$S$ をスキームとし、$(U, R, s, t, c)$ を $S$ 上の亜群スキームとする。
$U = \Spec(k)$ で $k$ が体であると仮定する。任意の点 $r \in R$ に対し、
\begin{enumerate}
\item $k'$ が代数閉である体拡大 $k'/k$、
\item $\Spec(k') \to \Spec(k)$ による $(U, R, s, t, c)$ の制限を
$(U', R', s', t', c')$ としたときの点 $r' \in R'$
\end{enumerate}
が存在し、次を満たす。
\begin{enumerate}
\item 点 $r'$ は射 $R' \to R$ により $r$ に写る。
\item 写像 $s', t' : R' \to \Spec(k')$ は同型 $k' \to \kappa(r')$ を誘導する。
\end{enumerate}
\end{lemma}

\begin{proof}
幾何学的主張を体についての主張に移すと、求めるものは図式
$$
\xymatrix{
k' & k' \ar[l]^1 & \\
k' \ar[u]^\tau & \kappa(r) \ar[lu]^\sigma & k \ar[l]^-s \ar[lu]_i \\
& k \ar[lu]^i \ar[u]_t
}
$$
である。ここで $i : k \to k'$ は $k$ の $k'$ への埋め込み、写像
$s, t : k \to \kappa(r)$ は $s, t : R \to U$ が誘導するもの、
$\tau : k' \to k'$ は自己同型である。この図式を次のように構成する。
\begin{enumerate}
\item $k'$ が代数閉で、$k$ 上の超越次数が十分大きいような体準同型
$i : k \to k'$ を選ぶ。
\item $\sigma \circ s = i$ を満たす埋め込み $\sigma : \kappa(r) \to k'$ を選ぶ。
このような $\sigma$ は次のように得られる。$\kappa(r)$ の $k$ 上の超越基底
$\{x_\alpha\}_{\alpha \in A}$ を選び、$i(k)$ 上代数的独立な元
$y_\alpha \in k'$, $\alpha \in A$ を取る。$\lambda \in k$ に対して
$s(\lambda) \mapsto i(\lambda)$、$\alpha \in A$ に対して
$x_\alpha \mapsto y_\alpha$ と定め、$s(k)(\{x_\alpha\})$ を $k'$ に写す。
次に $k'$ が代数閉であることを用いて $\tau : \kappa(\alpha) \to k'$ へ延長する。
\item $\tau \circ i = \sigma \circ t$ を満たす自己同型 $\tau : k' \to k'$ を選ぶ。
そのため、$k$ の素体上の超越基底 $\{x_\alpha\}_{\alpha \in A}$ を選ぶ。
一方で $\{i(x_\alpha)\}$ に $\{y_\beta\}_{\beta \in B}$ を加えて $k'$ の
超越基底へ延長し、他方で $\{\sigma(t(x_\alpha))\}$ に
$\{z_\gamma\}_{\gamma \in C}$ を加えて $k'$ の超越基底へ延長する。
$k'$ は代数閉なので、同型
$\sigma \circ t \circ i^{-1} : i(k) \to \sigma(t(k))$ は、$k'$ における
それぞれの代数閉包の同型
$\tau' : \overline{i(k)} \to \overline{\sigma(t(k))}$ へ延長できる。
$k'$ の超越次数は大きいので、集合 $B$ と $C$ の濃度は等しい。したがって全単射
$B \to C$ を用いて $\tau'$ を同型
$$
\overline{i(k)}(\{y_\beta\})
\longrightarrow
\overline{\sigma(t(k))}(\{z_\gamma\})
$$
へ延長できる。さらに $k'$ は両辺の代数閉包なので、これは望みどおり自己同型
$\tau : k' \to k'$ へ延長される。
\end{enumerate}
以上で補題が示された。
\end{proof}

\begin{lemma}
\label{more-groupoids-lemma-groupoid-on-field-move-point}
$S$ をスキームとし、$(U, R, s, t, c)$ を $S$ 上の亜群スキームとする。
$U = \Spec(k)$ で $k$ が体であると仮定する。$r \in R$ を、$s, t$ が同型
$k \to \kappa(r)$ を誘導する点とする。このとき写像
$$
R \longrightarrow R, \quad
x \longmapsto c(r, x)
$$
（精密な記法については証明を参照）は、$e$ を $r$ に写す自己同型 $R \to R$ である。
\end{lemma}

\begin{proof}
亜群を関手的に考えれば完全に明らかだが、すべて書き下す。
$a : U \to R$ を、像が $r$ で $s \circ a = \text{id}_U$ を満たす射とする。
これは $s : k \to \kappa(r)$ が同型であるという仮定により存在する。同様に、
$b : U \to R$ を、像が $r$ で $t \circ b = \text{id}_U$ を満たす射とする。
$b = a \circ (t \circ a)^{-1}$、特に $a \circ s \circ b = b$ である。

\medskip\noindent
$T$-値点上で
$$
(f : T \to R) \longmapsto (c(a \circ t \circ f, f) : T \to R)
$$
により与えられる射 $\Psi : R \to R$ を考える。これが定義されるには
$s \circ a \circ t \circ f = t \circ f$ を確認すればよいが、
$s \circ a = 1$ なので明らかである。$\Phi(e) = a$ であるから、補題を示すには
$\Phi$ が $R$ の自己同型であることを示せば十分である。
$T$-値点上で
$$
(g : T \to R) \longmapsto (c(i \circ b \circ t \circ g, g) : T \to R).
$$
により与えられる射を $\Phi : R \to R$ とする。これは
$s \circ i \circ b \circ t \circ g = t \circ b \circ t \circ g =
t \circ g$ なので定義される。$\Phi$ と $\Psi$ が互いに逆であることを示す。
実際、
\begin{align*}
& c(a \circ t \circ c(i \circ b \circ t \circ g, g),
c(i \circ b \circ t \circ g, g)) \\
& =
c(a \circ t \circ i \circ b \circ t \circ g,
c(i \circ b \circ t \circ g, g)) \\
& =
c(a \circ s \circ b \circ t \circ g,
c(i \circ b \circ t \circ g, g)) \\
& =
c(b \circ t \circ g, c(i \circ b \circ t \circ g, g)) \\
& =
c(c(b \circ t \circ g, i \circ b \circ t \circ g), g)) \\
& =
c(e, g) \\
& = g
\end{align*}
である。ここで上で示した関係 $a \circ s \circ b = b$ を用いた。逆向きには、
\begin{align*}
& c(i \circ b \circ t \circ c(a \circ t \circ f, f), c(a \circ t \circ f, f)) \\
& =
c(i \circ b \circ t \circ a \circ t \circ f, c(a \circ t \circ f, f)) \\
& =
c(i \circ a \circ (t \circ a)^{-1} \circ t \circ a \circ t \circ f,
c(a \circ t \circ f, f)) \\
& =
c(i \circ a \circ t \circ f, c(a \circ t \circ f, f)) \\
& =
c(c(i \circ a \circ t \circ f, a \circ t \circ f), f) \\
& =
c(e, f) \\
& = f
\end{align*}
を得る。以上で補題が示された。
\end{proof}

\begin{lemma}
\label{more-groupoids-lemma-groupoid-on-field-translate-open}
$S$ をスキームとし、$(U, R, s, t, c)$ を $S$ 上の亜群スキームとする。
$U$ が体のスペクトルで、$W \subset R$ が開、$Z \to R$ がスキームの射ならば、
合成 $Z \times_{s, U, t} W \to R \times_{s, U, t} R \to R$ の像は開である。
\end{lemma}

\begin{proof}
$U = \Spec(k)$ と書き、体拡大 $k'/k$ を考える。$U' = \Spec(k')$ と書き、
$R'$ を $U' \to U$ による $R$ の制限とする。
$Z' = Z \times_R R'$, $W' = R' \times_R W$ とおく。
$Z \times_{s, U, t} W$ の点 $\xi = (z, w)$ を取り、$Z \to R$ による $z$ の像を
$r \in R$ とする。補題 \ref{more-groupoids-lemma-groupoid-on-field-explain-points} のように
$k' \supset k$ と $r' \in R'$ を選ぶ。$z$ と $r'$ に写る $z' \in Z'$ を選べる。
すると $z'$ と $\xi$ に写る $\xi' \in Z' \times_{s', U', t'} W'$ を取れる。
開集合 $c(r', W')$（補題 \ref{more-groupoids-lemma-groupoid-on-field-move-point}）は
$Z' \times_{s', U', t'} W' \to R'$ の像に含まれる。ここで
$Z' \times_{s', U', t'} W' = (Z \times_{s, U, t} W)
\times_{R \times_{s, U, t} R} (R' \times_{s', U', t'} R')$ である。
したがって $Z' \times_{s', U', t'} W' \to R' \to R$ の像は
$Z \times_{s, U, t} W \to R$ の像に含まれる。$R' \to R$ は開なので
（補題 \ref{more-groupoids-lemma-restrict-groupoid-on-field}）、この像は望みどおり $\xi$ の像の
開近傍を含む。
\end{proof}

\begin{lemma}
\label{more-groupoids-lemma-groupoid-on-field-geometrically-irreducible}
$S$ をスキームとし、$(U, R, s, t, c)$ を $S$ 上の亜群スキームとする。
$U = \Spec(k)$ で $k$ が体であると仮定する。記法を濫用して、恒等射
$e : U \to R$ の像を $e \in R$ と書く。このとき、
\begin{enumerate}
\item $R$ の各局所環 $\mathcal{O}_{R, r}$ は一意な極小素イデアルをもち、
\item $e$ を通る $R$ の既約成分 $Z$ はちょうど一つであり、
\item $Z$ は $s$ または $t$ のいずれによっても $k$ 上幾何学的既約である。
\end{enumerate}
\end{lemma}

\begin{proof}
点 $r \in R$ を取る。本証明では、$r$ を通る $R$ の既約成分と局所環
$\mathcal{O}_{R, r}$ の極小素イデアルとの対応を、以後断りなく用いる。
補題 \ref{more-groupoids-lemma-groupoid-on-field-explain-points} のように
$k \subset k'$ と $r' \in R'$ を選ぶ。局所準同型
$\mathcal{O}_{R, r} \to \mathcal{O}_{R', r'}$ は忠実平坦である
（補題 \ref{more-groupoids-lemma-restrict-groupoid-on-field}）。したがって $r' \in R'$ についての
結果から $r \in R$ についての結果が従う。言い換えれば、
$s, t : k \to \kappa(r)$ が同型であると仮定してよい。補題
\ref{more-groupoids-lemma-groupoid-on-field-move-point} により、$e$ を $r$ に移す自己同型が存在する。
よって $r = e$ と仮定してよく、すなわち (1) は (2) から従う。

\medskip\noindent
まず $k$ が分離代数閉である場合に (2) を示す。$X, Y \subset R$ を $e$ を通る
既約成分とする。『多様体』補題
\ref{varieties-lemma-bijection-irreducible-components} および
\ref{varieties-lemma-separably-closed-irreducible} により、
$X \times_{s, U, t} Y$ も既約である。したがって
$c(X \times_{s, U, t} Y) \subset R$ は既約部分集合である。これは $X$ と $Y$ の
両方を $R$ の部分集合として含むと主張する。実際、$T$ を体のスペクトルとし、
$x : T \to X$ を $X$ の $T$-値点とする。このとき
$c(x, e \circ s \circ x) = x$ であり、$e \in Y$ なので
$e \circ s \circ x$ は $Y$ を経由する。$Y$ の点についても同様である。
よって明らかに $X = Y$、すなわち $e$ を通る $R$ の既約成分は一意である。

\medskip\noindent
一般の場合の (2), (3) を示す。$k \subset k'$ を分離代数閉包とし、
$(U', R', s', t', c')$ を $\Spec(k') \to \Spec(k)$ による
$(U, R, s, t, c)$ の制限とする。前段落により、$e'$ を通る $R'$ の既約成分
$Z'$ はちょうど一つである。$e$ の基底変換を
$e'' \in R \times_{s, U} U'$ と書く。補題
\ref{more-groupoids-lemma-restrict-groupoid-on-field} により
$R' \to R \times_{s, U} U'$ は忠実平坦で、$e' \mapsto e''$ なので、
$e''$ を通る $R \times_{s, k} k'$ の既約成分 $Z''$ はちょうど一つである。
$R \times_k k' \to R$ も忠実平坦だから、$e$ を通る $R$ の既約成分 $Z$ も
ちょうど一つである。これで (2) が示された。

\medskip\noindent
(3) を示すため、$Z''' \subset R \times_k k'$ を $Z \times_k k'$ の任意の既約成分とする。『多様体』補題
\ref{varieties-lemma-orbit-irreducible-components} により、$Z''' = \sigma(Z'')$ である。ここで $\sigma \in \text{Gal}(k'/k)$ である。
さらに $\sigma(e'') = e''$ なので、$e'' \in Z'''$ であり、したがって
$Z''' = Z''$ である。これは $Z$ が射 $s$ により $\Spec(k)$ 上幾何学的既約であることを
意味する。同じ議論により、$Z$ は射 $t$ によっても $\Spec(k)$ 上幾何学的既約である。
\end{proof}

\begin{lemma}
\label{more-groupoids-lemma-groupoid-on-field-locally-finite-type-dimension}
$S$ をスキームとし、$(U, R, s, t, c)$ を $S$ 上の亜群スキームとする。
$U = \Spec(k)$ で $k$ が体、$s, t$ が局所有限型であると仮定する。このとき、
\begin{enumerate}
\item $R$ は等次元である。
\item すべての $r \in R$ に対して $\dim(R) = \dim_r(R)$ である。
\item 任意の $r \in R$ に対して
$\text{trdeg}_{s(k)}(\kappa(r)) = \text{trdeg}_{t(k)}(\kappa(r))$ である。
\item 任意の閉点 $r \in R$ に対して
$\dim(R) = \dim(\mathcal{O}_{R, r})$ である。
\end{enumerate}
\end{lemma}

\begin{proof}
$r, r' \in R$ とする。補題 \ref{more-groupoids-lemma-groupoid-on-field-homogeneous} および
『射』補題 \ref{morphisms-lemma-dimension-fibre-after-base-change} により、
$\dim_r(R) = \dim_{r'}(R)$ である。『射』補題
\ref{morphisms-lemma-dimension-fibre-at-a-point} により
$$
\dim_r(R) =
\dim(\mathcal{O}_{R, r}) + \text{trdeg}_{s(k)}(\kappa(r)) =
\dim(\mathcal{O}_{R, r}) + \text{trdeg}_{t(k)}(\kappa(r)).
$$
他方、$R$（または $R$ の任意の開部分集合）の次元は、$R$ の局所環の次元の
上限である。『性質』補題 \ref{properties-lemma-codimension-local-ring} を参照せよ。
これは明らかに閉点 $r$ で最大となり、その場合 Hilbert の零点定理により
$\text{trdeg}_k(\kappa(r)) = 0$ である（『射』節
\ref{morphisms-section-points-finite-type} を参照）。よって補題が従う。
\end{proof}

\begin{lemma}
\label{more-groupoids-lemma-groupoid-on-field-dimension-equal-stabilizer}
$S$ をスキームとし、$(U, R, s, t, c)$ を $S$ 上の亜群スキームとする。
$U = \Spec(k)$ で $k$ が体、$s, t$ が局所有限型であると仮定する。
$G$ を $R$ の安定化群スキームとすると、$\dim(R) = \dim(G)$ である。
\end{lemma}

\begin{proof}
$Z \subset R$ を $e$ を通る既約成分とし（補題
\ref{more-groupoids-lemma-groupoid-on-field-geometrically-irreducible}）、$R$ の整閉部分スキームと
みなす。$k'_s$, それぞれ\ $k'_t$ を、$\Gamma(Z, \mathcal{O}_Z)$ における
$s(k)$, それぞれ\ $t(k)$ の整閉包とする。$k'_s$ と $k'_t$ は体である。
『多様体』補題 \ref{varieties-lemma-integral-closure-ground-field} を参照せよ。
『多様体』命題 \ref{varieties-proposition-unique-base-field} により、
$\Gamma(Z, \mathcal{O}_Z)$ の部分環として $k'_s = k'_t$ である。
$e$ は $Z$ を経由するので、可換図式
$$
\xymatrix{
k \ar[rd]_t \ar[rrd]^1 \\
& \Gamma(Z, \mathcal{O}_Z) \ar[r]^e & k \\
k \ar[ru]^s \ar[rru]_1
}
$$
を得る。一方でこれは $k'_s = s(k)$, $k'_t = t(k)$、したがって
$s(k) = t(k)$ を示す。これを上の図式と合わせると $s = t$ である！
言い換えれば、$Z$ は
$G = R \times_{(t, s), U \times_S U, \Delta} U$ の閉部分スキームである。
$G$ と $R$ はとも等次元なので補題が従う。補題
\ref{more-groupoids-lemma-groupoid-on-field-locally-finite-type-dimension} および
『亜群スキーム』補題 \ref{groupoids-lemma-group-scheme-finite-type-field} を参照せよ。
\end{proof}

\begin{remark}
\label{more-groupoids-remark-warn-dimension-groupoid-on-field}
注意：補題 \ref{more-groupoids-lemma-groupoid-on-field-dimension-equal-stabilizer} は、$s$ と $t$ が
局所有限型であるという条件なしには誤りである。簡単な例として、作用
$$
\mathbf{G}_{m, \mathbf{Q}} \times_{\mathbf{Q}} \mathbf{A}^1_{\mathbf{Q}}
\to \mathbf{A}^1_{\mathbf{Q}}
$$
から始め、対応する亜群スキームを $\mathbf{A}^1_{\mathbf{Q}}$ の一般点に制限する。
言い換えれば、射
$\Spec(\mathbf{Q}(x)) \to
\Spec(\mathbf{Q}[x]) = \mathbf{A}^1_{\mathbf{Q}}$
により制限する。すると $U = \Spec(\mathbf{Q}(x))$ および
$$
R = \Spec\left(
\mathbf{Q}(x)[y]\left[
\frac{1}{P(xy)}, P \in \mathbf{Q}[T], P \not = 0
\right]
\right)
$$
を満たす亜群スキーム $(U, R, s, t, c)$ を得る。この場合
$\dim(R) = 1$, $\dim(G) = 0$ である。
\end{remark}

\begin{lemma}
\label{more-groupoids-lemma-groupoid-characteristic-zero-smooth}
$S$ をスキームとし、$(U, R, s, t, c)$ を $S$ 上の亜群スキームとする。
次を仮定する。
\begin{enumerate}
\item $U = \Spec(k)$ で $k$ は体である。
\item $s, t$ は局所有限型である。
\item $k$ の標数は零である。
\end{enumerate}
このとき $s, t : R \to U$ は滑らかである。
\end{lemma}

\begin{proof}
補題 \ref{more-groupoids-lemma-sheaf-differentials} により、$R \to U$ の微分層は自由である。
したがって『多様体』補題
\ref{varieties-lemma-char-zero-differentials-free-smooth} から滑らかさが従う。
\end{proof}

\begin{lemma}
\label{more-groupoids-lemma-reduced-group-scheme-perfect-field-characteristic-p-smooth}
$S$ をスキームとし、$(U, R, s, t, c)$ を $S$ 上の亜群スキームとする。
次を仮定する。
\begin{enumerate}
\item $U = \Spec(k)$ で $k$ は体である。
\item $s, t$ は局所有限型である。
\item $R$ は被約である。
\item $k$ は完全である。
\end{enumerate}
このとき $s, t : R \to U$ は滑らかである。
\end{lemma}

\begin{proof}
補題 \ref{more-groupoids-lemma-sheaf-differentials} により、層 $\Omega_{R/U}$ は自由である。
したがって『多様体』補題
\ref{varieties-lemma-char-p-differentials-free-smooth} から補題が従う。
\end{proof}










\section{体上の亜群の射}
\label{more-groupoids-section-morphisms-groupoids-on-fields}

\noindent
本節では体上の亜群の間の射を調べる。これは体上の群スキームの射を調べる場合より
少し一般的だが、非常によく似ている。

\begin{situation}
\label{more-groupoids-situation-morphism-groupoids-on-field}
$S$ をスキームとし、$U = \Spec(k)$ を $S$ 上のスキームとする。ここで $k$ は体である。
$(U, R_1, s_1, t_1, c_1)$, $(U, R_2, s_2, t_2, c_2)$ を、第一成分が同じ
$S$ 上の亜群スキームとする。$a : R_1 \to R_2$ を、$(\text{id}_U, a)$ が
$S$ 上の亜群スキームの射を定めるような射とする。『亜群スキーム』定義
\ref{groupoids-definition-groupoid} を参照せよ。特に次の図式は可換である。
$$
\vcenter{
\xymatrix{
R_1 \ar[rrd]^{t_1} \ar[rdd]_{s_1} \ar[rd]_a \\
& R_2 \ar[d]^{t_2} \ar[r]_{s_2} & U \\
& U
}
}
\quad\quad
\vcenter{
\xymatrix{
R_1 \times_{s_1, U, t_1} R_1 \ar[r]_-{c_1} \ar[d]_{a \times a} &
R_1 \ar[d]^a \\
R_2 \times_{s_2, U, t_2} R_2 \ar[r]^-{c_2} &
R_2
}
}
$$
\end{situation}

\noindent
次の補題は『亜群スキーム』補題
\ref{groupoids-lemma-open-subgroup-closed-over-field} の一般化である。

\begin{lemma}
\label{more-groupoids-lemma-open-image-is-closed}
記法と仮定は状況 \ref{more-groupoids-situation-morphism-groupoids-on-field} のとおりとする。
$a(R_1)$ が $R_2$ で開ならば、$a(R_1)$ は $R_2$ で閉である。
\end{lemma}

\begin{proof}
$r_2 \in R_2$ を $a(R_1)$ の閉包に属する点とする。$r_2 \in a(R_1)$ を示したい。
補題 \ref{more-groupoids-lemma-groupoid-on-field-explain-points} のように、
$(U, R_2, s_2, t_2, c_2)$ と $r_2$ に適合する $k \subset k'$ および
$r_2' \in R'_2$ を選ぶ。$R_i'$ を射
$U' = \Spec(k') \to U = \Spec(k)$ による $R_i$ の制限とする。
$a' : R'_1 \to R_2'$ を $a$ の基底変換とする。図式
$$
\xymatrix{
R'_1 \ar[r]_{a'} \ar[d]_{p_1} & R'_2 \ar[d]^{p_2} \\
R_1 \ar[r]^a & R_2
}
$$
はファイバー正方形である。したがって $a'$ の像は、射 $p_2 : R'_2 \to R_2$ による
$a$ の像の逆像である。補題 \ref{more-groupoids-lemma-restrict-groupoid-on-field} により $p_2$ は
全射かつ開である。よって『位相』補題
\ref{topology-lemma-open-morphism-quotient-topology} により、$r_2'$ は
$a'(R'_1)$ の閉包に属する。したがって、$r_2 \in R_2$ について、$s_2$ と $t_2$ が誘導する写像
$k \to \kappa(r_2)$ が同型であると仮定してよい。

\medskip\noindent
この場合、補題 \ref{more-groupoids-lemma-groupoid-on-field-move-point} を用いられる。この補題から
$c(r_2, a(R_1))$ は $r_2$ の開近傍である。$r_2$ は $a(R_1)$ の閉包の点だから、
$a(R_1) \cap c(r_2, a(R_1)) \not = \emptyset$ である。$R_2$ と $R_1$ の逆元写像を
用いると、これは $c_2(a(R_1), a(R_1))$ が $r_2$ を含むことを意味する。
$c_2(a(R_1), a(R_1)) \subset a(c_1(R_1, R_1)) = a(R_1)$ なので、望みどおり
$r_2 \in a(R_1)$ を得る。
\end{proof}

\begin{lemma}
\label{more-groupoids-lemma-map-groupoids-on-field-image}
記法と仮定は状況 \ref{more-groupoids-situation-morphism-groupoids-on-field} のとおりとする。
$Z \subset R_2$ を、台位相空間が $a : R_1 \to R_2$ の像の閉包であるような被約閉部分
スキームとする（『スキーム』定義 \ref{schemes-definition-reduced-induced-scheme}）。
このとき集合論的に $c_2(Z \times_{s_2, U, t_2} Z) \subset Z$ である。
\end{lemma}

\begin{proof}
可換図式
$$
\xymatrix{
R_1 \times_{s_1, U, t_1} R_1 \ar[r] \ar[d] & R_1 \ar[d] \\
R_2 \times_{s_2, U, t_2} R_2 \ar[r] & R_2
}
$$
を考える。『多様体』補題 \ref{varieties-lemma-closure-image-product-map} により、
左の垂直射の像の閉包は集合論的に $Z \times_{s_2, U, t_2} Z$ である。
よって結果が従う。
\end{proof}

\begin{lemma}
\label{more-groupoids-lemma-map-groupoids-on-perfect-field-image}
記法と仮定は状況 \ref{more-groupoids-situation-morphism-groupoids-on-field} のとおりとし、$k$ は完全とする。
$Z \subset R_2$ を、台位相空間が $a : R_1 \to R_2$ の像の閉包であるような被約閉部分
スキームとする（『スキーム』定義 \ref{schemes-definition-reduced-induced-scheme}）。
このとき
$$
(U, Z, s_2|_Z, t_2|_Z, c_2|_Z)
$$
は $S$ 上の亜群スキームである。
\end{lemma}

\begin{proof}
まず主張が意味をもつことを説明する。$U$ は完全体 $k$ のスペクトルなので、
$Z$ はいずれの射影によっても $k$ 上幾何学的被約である。『多様体』補題
\ref{varieties-lemma-perfect-reduced} を参照せよ。したがってスキーム
$Z \times_{s_2, U, t_2} Z \subset Z$ は被約である。『多様体』補題
\ref{varieties-lemma-geometrically-reduced-any-base-change} を参照せよ。
よって補題 \ref{more-groupoids-lemma-map-groupoids-on-field-image} により、$c$ は射
$Z \times_{s_2, U, t_2} Z \to Z$ を誘導する。最後に、$e_2$ が $Z$ を経由し、
写像 $i_2 : R_2 \to R_2$ が $Z$ を保つことは明らかである。七つ組
$(U, R_2, s_2, t_2, c_2, e_2, i_2)$ の射は亜群の公理を満たすので、$Z$ に制限した後も
公理を満たす。
\end{proof}

\begin{lemma}
\label{more-groupoids-lemma-locally-closed-image-is-closed}
記法と仮定は状況 \ref{more-groupoids-situation-morphism-groupoids-on-field} のとおりとする。
像 $a(R_1)$ が $R_2$ の局所閉部分集合ならば、それは閉部分集合である。
\end{lemma}

\begin{proof}
$k \subset k'$ を体 $k$ の完全閉包とする。$R_i'$ を射
$U' = \Spec(k') \to \Spec(k)$ による $R_i$ の制限とする。射 $R_i' \to R_i$ は、
普遍的同相射 $U' \to U$ の基底変換の合成なので普遍的同相射である
（補題 \ref{more-groupoids-lemma-restrict-groupoid-on-field} の主張中の図式を参照）。
したがって $a'(R_1')$ が $R_2'$ で閉であることを示せば十分である。
言い換えれば、$k$ が完全であると仮定してよい。

\medskip\noindent
$k$ が完全ならば、補題 \ref{more-groupoids-lemma-map-groupoids-on-perfect-field-image} により、
像の閉包は亜群スキーム $Z \subset R_2$ である。同じ補題を
$\text{id}_{R_1} : R_1 \to R_1$ に適用すると、$(R_2)_{red}$ は亜群スキームである。
そこで射 $a|_{(R_2)_{red}} : (R_2)_{red} \to Z$ に補題
\ref{more-groupoids-lemma-open-image-is-closed} を適用し、$Z$ が $a$ の像に等しいと結論する。
\end{proof}

\begin{lemma}
\label{more-groupoids-lemma-quasi-compact-map-groupoids-on-field-image}
記法と仮定は状況 \ref{more-groupoids-situation-morphism-groupoids-on-field} のとおりとする。
$a : R_1 \to R_2$ が準コンパクト射であると仮定する。
$Z \subset R_2$ を $a : R_1 \to R_2$ のスキーム論的像とする
（『射』定義 \ref{morphisms-definition-scheme-theoretic-image} を参照）。このとき
$$
(U, Z, s_2|_Z, t_2|_Z, c_2|_Z)
$$
は $S$ 上の亜群スキームである。
\end{lemma}

\begin{proof}
主な難点は、$c_2|_{Z \times_{s_2, U, t_2} Z}$ が $Z$ に入ることを示す点にある。
可換図式
$$
\xymatrix{
R_1 \times_{s_1, U, t_1} R_1 \ar[r] \ar[d]^{a \times a} & R_1 \ar[d] \\
R_2 \times_{s_2, U, t_2} R_2 \ar[r] & R_2
}
$$
を考える。『多様体』補題
\ref{varieties-lemma-scheme-theoretic-image-product-map} により、
$a \times a$ のスキーム論的像は $Z \times_{s_2, U, t_2} Z$ である。
図式の可換性から、下の水平射は $Z \times_{s_2, U, t_2} Z$ を $Z$ に写す。
補題 \ref{more-groupoids-lemma-map-groupoids-on-perfect-field-image} の証明と同様に、
$i_2(Z) \subset Z$ であり、かつ $e_2$ は $Z$ を経由する。したがって同補題の証明と
同様に結論を得る。
\end{proof}

\begin{lemma}
\label{more-groupoids-lemma-groupoid-on-field-image}
$S$ をスキームとし、$(U, R, s, t, c)$ を $S$ 上の亜群スキームとする。
$U$ は体のスペクトルであると仮定する。$Z \subset U \times_S U$ を、台位相空間が
$j = (t, s) : R \to U \times_S U$ の像の閉包であるような被約閉部分スキームとする
（『スキーム』定義 \ref{schemes-definition-reduced-induced-scheme} を参照）。このとき集合論的に
$\text{pr}_{02}(Z \times_{\text{pr}_1, U, \text{pr}_0} Z) \subset Z$
である。
\end{lemma}

\begin{proof}
$(U, U \times_S U, \text{pr}_1, \text{pr}_0, \text{pr}_{02})$ は $S$ 上の
亜群スキームなので、これは補題 \ref{more-groupoids-lemma-map-groupoids-on-field-image} の特別な場合である。
次のように直接証明することもできる。

\medskip\noindent
$U = \Spec(k)$ と書く。$R_s$（resp.\ $Z_s$，resp.\ $U^2_s$）で、
$R$（resp.\ $Z$，resp.\ $U \times_S U$）を $s$（resp.\ $\text{pr}_1|_Z$，
resp.\ $\text{pr}_1$）を通じて $k$ 上のスキームとみなしたものを表す。同様に、
${}_tR$（resp.\ ${}_tZ$，resp.\ ${}_tU^2$）で、
$R$（resp.\ $Z$，resp.\ $U \times_S U$）を $t$（resp.\ $\text{pr}_0|_Z$，
resp.\ $\text{pr}_0$）を通じて $k$ 上のスキームとみなしたものを表す。
射 $j$ は $k$ 上のスキームの射
$j_s : R_s \to U^2_s$ および ${}_tj : {}_tR \to {}_tU^2$ を誘導する。
可換図式
$$
\xymatrix{
R_s \times_k {}_tR \ar[r]^c \ar[d]_{j_s \times {}_tj} &  R \ar[d]^j \\
U^2_s \times_k {}_tU^2 \ar[r] & U \times_S U
}
$$
を考える。『多様体』補題 \ref{varieties-lemma-closure-image-product-map} により、
$j_s \times {}_tj$ の像の閉包は $Z_s \times_k {}_tZ$ である。図式の可換性から、
下の水平射は $Z_s \times_k {}_tZ$ を $Z$ に写す。
\end{proof}

\begin{lemma}
\label{more-groupoids-lemma-groupoid-on-perfect-field-image}
$S$ をスキームとし、$(U, R, s, t, c)$ を $S$ 上の亜群スキームとする。
$U$ は完全体のスペクトルであると仮定する。$Z \subset U \times_S U$ を、台位相空間が
$j = (t, s) : R \to U \times_S U$ の像の閉包であるような被約閉部分スキームとする
（『スキーム』定義 \ref{schemes-definition-reduced-induced-scheme} を参照）。このとき
$$
(U, Z, \text{pr}_0|_Z, \text{pr}_1|_Z,
\text{pr}_{02}|_{Z \times_{\text{pr}_1, U, \text{pr}_0} Z})
$$
は $S$ 上の亜群スキームである。
\end{lemma}

\begin{proof}
$(U, U \times_S U, \text{pr}_1, \text{pr}_0, \text{pr}_{02})$ は $S$ 上の
亜群スキームなので、これは補題 \ref{more-groupoids-lemma-map-groupoids-on-perfect-field-image} の特別な場合である。
次のように直接証明することもできる。

\medskip\noindent
まず主張が意味をもつことを説明する。$U$ は完全体 $k$ のスペクトルなので、スキーム $Z$ は
いずれの射影によっても $k$ 上幾何学的被約である。『多様体』補題
\ref{varieties-lemma-perfect-reduced} を参照せよ。したがってスキーム
$Z \times_{\text{pr}_1, U, \text{pr}_0} Z \subset Z$ は被約である。『多様体』補題
\ref{varieties-lemma-geometrically-reduced-any-base-change} を参照せよ。
よって補題 \ref{more-groupoids-lemma-groupoid-on-field-image} により、$\text{pr}_{02}$ は射
$Z \times_{\text{pr}_1, U, \text{pr}_0} Z \to Z$ を誘導する。最後に、
$\Delta_{U/S}$ が $Z$ を経由し、写像
$\sigma : U \times_S U \to U \times_S U$, $(x, y) \mapsto (y, x)$
が $Z$ を保つことは明らかである。七つ組
$(U, U \times_S U, \text{pr}_0, \text{pr}_1, \text{pr}_{02},
\Delta_{U/S}, \sigma)$ は亜群の公理を満たすので、$Z$ に制限した後も公理を満たす。
\end{proof}

\begin{lemma}
\label{more-groupoids-lemma-quasi-compact-groupoid-on-field-image}
$S$ をスキームとし、$(U, R, s, t, c)$ を $S$ 上の亜群スキームとする。
$U$ は体のスペクトルであり、$R$ は準コンパクトであると仮定する
（これは $s, t$ が準コンパクトであることと同値である）。$Z \subset U \times_S U$ を
$j = (t, s) : R \to U \times_S U$ のスキーム論的像とする
（『射』定義 \ref{morphisms-definition-scheme-theoretic-image} を参照）。このとき
$$
(U, Z, \text{pr}_0|_Z, \text{pr}_1|_Z,
\text{pr}_{02}|_{Z \times_{\text{pr}_1, U, \text{pr}_0} Z})
$$
は $S$ 上の亜群スキームである。
\end{lemma}

\begin{proof}
$(U, U \times_S U, \text{pr}_1, \text{pr}_0, \text{pr}_{02})$ は $S$ 上の
亜群スキームなので、これは補題 \ref{more-groupoids-lemma-quasi-compact-map-groupoids-on-field-image} の特別な場合である。
次のように直接証明することもできる。

\medskip\noindent
主な難点は、$\text{pr}_{02}|_{Z \times_{\text{pr}_1, U, \text{pr}_0} Z}$ が
$Z$ に入ることを示す点にある。$U = \Spec(k)$ と書く。
$R_s$（resp.\ $Z_s$，resp.\ $U^2_s$）で、
$R$（resp.\ $Z$，resp.\ $U \times_S U$）を $s$（resp.\ $\text{pr}_1|_Z$，
resp.\ $\text{pr}_1$）を通じて $k$ 上のスキームとみなしたものを表す。同様に、
${}_tR$（resp.\ ${}_tZ$，resp.\ ${}_tU^2$）で、
$R$（resp.\ $Z$，resp.\ $U \times_S U$）を $t$（resp.\ $\text{pr}_0|_Z$，
resp.\ $\text{pr}_0$）を通じて $k$ 上のスキームとみなしたものを表す。
射 $j$ は $k$ 上のスキームの射
$j_s : R_s \to U^2_s$ および ${}_tj : {}_tR \to {}_tU^2$ を誘導する。
可換図式
$$
\xymatrix{
R_s \times_k {}_tR \ar[r]^c \ar[d]_{j_s \times {}_tj} &  R \ar[d]^j \\
U^2_s \times_k {}_tU^2 \ar[r] & U \times_S U
}
$$
を考える。『多様体』補題 \ref{varieties-lemma-scheme-theoretic-image-product-map} により、
$j_s \times {}_tj$ のスキーム論的像は $Z_s \times_k {}_tZ$ である。図式の可換性から、
下の水平射は $Z_s \times_k {}_tZ$ を $Z$ に写す。補題
\ref{more-groupoids-lemma-groupoid-on-perfect-field-image} の証明と同様に、$\sigma(Z) \subset Z$ であり、
かつ $\Delta_{U/S}$ は $Z$ を経由する。したがって同補題の証明と同様に結論を得る。
\end{proof}







\section{亜群のスライス}
\label{more-groupoids-section-slicing}

\noindent
次の補題は、安定化群の次元が小さいとき、Cohen--Macaulay 亜群スキームをスライスして
ファイバーの次元を下げられることを示す。これは、商スタックの与えられた表示を改善する
過程における本質的な一歩である。

\begin{situation}
\label{more-groupoids-situation-slice}
$S$ をスキームとし、$(U, R, s, t, c)$ を $S$ 上の亜群スキームとする。
$g : U' \to U$ をスキームの射とする。$u \in U$ を点とし、$u' \in U'$ を
$g(u') = u$ を満たす点とする。これらのデータに対し、$(U', R', s', t', c')$ で
射 $g$ による $(U, R, s, t, c)$ の制限を表す。$G \to U$ で $R$ の安定化群
スキームを表す。これは $R$ の局所閉部分スキームである。$h$ で合成
$$
h = s \circ \text{pr}_1 : U' \times_{g, U, t} R \longrightarrow U.
$$
を表す。$F_u = s^{-1}(u)$（スキーム論的ファイバー）とおき、$G_u$ で $u$ 上の
$G$ のスキーム論的ファイバーを表す。同様に、$R'$ に対して
$F'_{u'} = (s')^{-1}(u')$ とおく。$g(u') = u$ なので
$$
F'_{u'} = h^{-1}(u) \times_{\Spec(\kappa(u))} \Spec(\kappa(u')).
$$
である。点 $e(u) \in R$ は $G_u$ および $F_u$ の点ともみなせる。また $e'(u')$ は
$R'$（resp.\ $G'_{u'}$，resp.\ $F'_{u'}$）の点であり、$R$（resp.\ $G_u$，
resp.\ $F_u$）において $e(u)$ に写る。
\end{situation}

\begin{lemma}
\label{more-groupoids-lemma-slice}
$S$ をスキームとし、$(U, R, s, t, c, e, i)$ を $S$ 上の亜群スキームとする。
$G \to U$ を安定化群スキームとする。$s$ と $t$ は Cohen--Macaulay かつ局所有限表示で
あると仮定する。$u \in U$ をスキーム $U$ の有限型点とする。『射』定義
\ref{morphisms-definition-finite-type-point} を参照せよ。状況 \ref{more-groupoids-situation-slice} の記法で
$$
d_1 = \dim(G_u), \quad
d_2 = \dim_{e(u)}(F_u).
$$
とおく。$d_2 > d_1$ ならば、アフィンスキーム $U'$ と射 $g : U' \to U$ であって
（記法は状況 \ref{more-groupoids-situation-slice} のとおり）次を満たすものが存在する。
\begin{enumerate}
\item $g$ は埋め込みである。
\item $u \in U'$ である。
\item $g$ は局所有限表示である。
\item 射 $h : U' \times_{g, U, t} R \longrightarrow U$ は $(u, e(u))$ で
Cohen--Macaulay である。
\item $\dim_{e'(u)}(F'_u) = d_2 - 1$ である。
\end{enumerate}
\end{lemma}

\begin{proof}
$\Spec(A) \subset U$ を $u$ のアフィン近傍で、$u$ が $U$ の閉点に対応するものとする。
『射』補題 \ref{morphisms-lemma-identify-finite-type-points} を参照せよ。
$\Spec(B) \subset R$ を $e(u)$ のアフィン近傍で、$j$ により開集合
$\Spec(A) \times_S \Spec(A) \subset U \times_S U$ に写るものとする。
$\mathfrak m \subset A$ を $u$ に対応する極大イデアル、$\mathfrak q \subset B$ を
$e(u)$ に対応する素イデアルとする。図式で書けば
$$
\vcenter{
\xymatrix{
B & A \ar[l]^s \\
A \ar[u]^t
}
}
\quad\text{and}\quad
\vcenter{
\xymatrix{
B_{\mathfrak q} & A_{\mathfrak m} \ar[l]^s \\
A_{\mathfrak m} \ar[u]^t
}
}
$$
である。二つの誘導写像
$s, t : \kappa(\mathfrak m) \to \kappa(\mathfrak q)$
は等しく、$s \circ e = t \circ e = \text{id}_U$ なので同型である。特に
$\mathfrak q$ も極大イデアルである。環準同型 $s, t : A \to B$ は有限表示かつ平坦である。
仮定により環
$$
\mathcal{O}_{F_u, e(u)} = B_{\mathfrak q}/s(\mathfrak m)B_{\mathfrak q}
$$
は次元 $d_2$ の Cohen--Macaulay 環である。次元の等式は『射』補題
\ref{morphisms-lemma-dimension-fibre-at-a-point} による。

\medskip\noindent
$R''$ を、射 $\Spec(\kappa(u)) \to U$ による $u = \Spec(\kappa(u))$ への $R$ の
制限とする。$u \to U$ は局所有限型なので、補題 \ref{more-groupoids-lemma-restrict-preserves-type} により
$(\Spec(\kappa(u)), R'', s'', t'', c'')$ は $s'', t''$ が局所有限型である亜群スキームである。
補題 \ref{more-groupoids-lemma-groupoid-on-field-dimension-equal-stabilizer} により
$\dim(G'') = \dim(R'')$ である。また補題
\ref{more-groupoids-lemma-groupoid-on-field-locally-finite-type-dimension} により
$\dim(R'') = \dim_{e''}(R'') = \dim(\mathcal{O}_{R'', e''})$ である。
『亜群スキーム』補題 \ref{groupoids-lemma-restrict-stabilizer} により $G'' = G_u$ である。
したがって $\dim(\mathcal{O}_{R'', e''}) = d_1$ を得る。

\medskip\noindent
スキームとして $R''$ は
$$
R'' =
R \times_{(U \times_S U)}
\Big(
\Spec(\kappa(\mathfrak m)) \times_S \Spec(\kappa(\mathfrak m))
\Big)
$$
である。したがって $e''$ のあるアフィン開近傍は環
$$
B \otimes_{(A \otimes A)} (\kappa(\mathfrak m) \otimes \kappa(\mathfrak m))
=
B/s(\mathfrak m)B + t(\mathfrak m)B
$$
のスペクトルである。ゆえに
$$
\mathcal{O}_{R'', e''} =
B_{\mathfrak q}/s(\mathfrak m)B_{\mathfrak q} + t(\mathfrak m)B_{\mathfrak q}
$$
であり、この環の次元は $d_1$ である。

\medskip\noindent
これより、次を満たす元 $f \in \mathfrak m$ が存在すると主張する。
$$
\dim(B_{\mathfrak q}/(s(\mathfrak m)B_{\mathfrak q} + fB_{\mathfrak q}) < d_2
$$
実際、$\mathfrak n_j \supset s(\mathfrak m)B_{\mathfrak q}$（$j = 1, \ldots, m$）が
局所環 $B_{\mathfrak q}/s(\mathfrak m)B_{\mathfrak q}$ の極小素イデアルに対応するとする。
この環は Noether 環なので、その個数は有限である（この環は体上本質的有限型であり、また
Cohen--Macaulay 環は Noether 環でもある）。Cohen--Macaulay 条件と『代数』補題
\ref{algebra-lemma-CM-dim-formula} により
$\dim(B_{\mathfrak q}/\mathfrak n_j) = d_2$ である。また
$\dim(B_{\mathfrak q}/(\mathfrak n_j + t(\mathfrak m)B_{\mathfrak q}))
\leq d_1$
である。実際、これは次元 $d_1$ の環
$\mathcal{O}_{R'', e''} =
B_{\mathfrak q}/s(\mathfrak m)B_{\mathfrak q} + t(\mathfrak m)B_{\mathfrak q}$
の商である。$d_1 < d_2$ なので $\mathfrak m \not \subset t^{-1}(\mathfrak n_i)$ である。
素回避定理（『代数』補題 \ref{algebra-lemma-silly}）により、すべての
$j = 1, \ldots, m$ に対して $t(f) \not \in \mathfrak n_j$ を満たす
$f \in \mathfrak m$ を選べる。この $f$ に対して上の不等式が成り立つ。
『代数』補題 \ref{algebra-lemma-one-equation} を参照せよ。

\medskip\noindent
$A' = A/fA$ および $U' = \Spec(A')$ とおく。このとき $U' \to U$ は埋め込みで
局所有限表示であり、$u \in U'$ であることは明らかである。したがって補題の (1), (2), (3) が成り立つ。
射
$$
U' \times_{g, U, t} R \longrightarrow U
$$
は $\Spec(A)$ を経由し、環準同型
$$
\xymatrix{
B/t(f)B \ar@{=}[r] & A/(f) \otimes_{A, t} B & A \ar[l]_-s
}
$$
に対応する。ここで、$B_{\mathfrak q}/s(\mathfrak m)B_{\mathfrak q}$ は正次元の
Cohen--Macaulay 環であり、$f$ はどの極小素イデアルにも含まれないので、$t(f)$ は
この環上の零因子ではない。例えば『代数』補題 \ref{algebra-lemma-reformulate-CM} を参照せよ。
したがって『代数』補題 \ref{algebra-lemma-grothendieck-general} により、
$s : A_{\mathfrak m} \to B_{\mathfrak q}/t(f)B_{\mathfrak q}$ は平坦であり、そのファイバー環は
$B_{\mathfrak q}/(s(\mathfrak m)B_{\mathfrak q} + t(f)B_{\mathfrak q})$
である。この環は再び『代数』補題 \ref{algebra-lemma-reformulate-CM} により
Cohen--Macaulay である。これで補題の (4) が従う。(5) を示すには、図式
(\ref{more-groupoids-equation-restriction}) によりファイバー $F'_u$ が $u$ 上の $h$ のファイバーに
等しいことに注意する。したがって『射』補題
\ref{morphisms-lemma-dimension-fibre-at-a-point} により
$\dim_{e'(u)}(F'_u) =
\dim(B_{\mathfrak q}/(s(\mathfrak m)B_{\mathfrak q} + t(f)B_{\mathfrak q}))$
であり、この環の次元は、もう一度『代数』補題
\ref{algebra-lemma-reformulate-CM} を用いると $d_2 - 1$ である。これで補題の最後の主張も
示された。
\end{proof}

\noindent
スライスの方法が分かったので、これを先の結果と組み合わせて次の「最適」な結果を得る。
ここで最適というのは、$G_u$ は $F_u$ の局所閉部分スキームなので常に
$\dim(G_u) = \dim_{e(u)}(G_u) \leq \dim_{e(u)}(F_u)$ であり、補題よりさらに
スライスすることはできない、という意味である。

\begin{lemma}
\label{more-groupoids-lemma-max-slice}
$S$ をスキームとし、$(U, R, s, t, c, e, i)$ を $S$ 上の亜群スキームとする。
$G \to U$ を安定化群スキームとする。$s$ と $t$ は Cohen--Macaulay かつ局所有限表示で
あると仮定する。$u \in U$ をスキーム $U$ の有限型点とする。『射』定義
\ref{morphisms-definition-finite-type-point} を参照せよ。状況 \ref{more-groupoids-situation-slice} の記法で、
次を満たすアフィンスキーム $U'$ と射 $g : U' \to U$ が存在する。
\begin{enumerate}
\item $g$ は埋め込みである。
\item $u \in U'$ である。
\item $g$ は局所有限表示である。
\item 射 $h : U' \times_{g, U, t} R \longrightarrow U$ は Cohen--Macaulay かつ
局所有限表示である。
\item 射 $s', t' : R' \to U'$ は Cohen--Macaulay かつ局所有限表示である。
\item $\dim_{e(u)}(F'_u) = \dim(G'_u)$ である。
\end{enumerate}
\end{lemma}

\begin{proof}
$s$ は局所有限表示なので、スキーム $F_u$ は $\kappa(u)$ 上局所有限型である。したがって
$\dim_{e(u)}(F_u) < \infty$ であり、$\dim_{e(u)}(F_u)$ に関する帰納法を用いられる。

\medskip\noindent
$\dim_{e(u)}(F_u) = \dim(G_u)$ なら示すことはない。
$\dim_{e(u)}(F_u) > \dim(G_u)$ と仮定する。このとき補題 \ref{more-groupoids-lemma-slice} が適用でき、
(1), (2), (3) を満たす射 $g : U' \to U$ が得られる。(6) の代わりに
$\dim_{e(u)}(F'_u) < \dim_{e(u)}(F_u)$ が成り立ち、(4), (5) の代わりに合成
$$
h = s \circ \text{pr}_1 : U' \times_{g, U, t} R \longrightarrow U
$$
が点 $(u, e(u))$ で Cohen--Macaulay となる。注意 \ref{more-groupoids-remark-local-source-apply} を
適用すると、$U'' \subset U'$ であって、
$U'' \times_{g, U, t} R \subset U' \times_{g, U, t} R$ が $h$ の Cohen--Macaulay
軌跡である最大の開部分スキームとなるものを得る。
$(u, e(u)) \in U'' \times_{g, U, t} R$ なので $u \in U''$ である。
したがって $U'$ を $U''$ で置き換え、$h$ は実際に至る所 Cohen--Macaulay であると
仮定してよい。補題 \ref{more-groupoids-lemma-restrict-property} より、$s', t'$ は局所有限表示かつ
Cohen--Macaulay である（『射』補題
\ref{morphisms-lemma-base-change-finite-presentation} および『射の続論』補題
\ref{more-morphisms-lemma-base-change-CM} を用いる）。

\medskip\noindent
構成により $\dim_{e'(u)}(F'_u) < \dim_{e(u)}(F_u)$ なので、
$(U', R', s', t', c')$ と点 $u \in U'$ に帰納法の仮定を適用できる。なお $u$ は
$U'$ の有限型点でもある（例えば『射』補題
\ref{morphisms-lemma-identify-finite-type-points} による有限型点の特徴付けを用いればよい）。
$g' : U'' \to U'$ と $(U'', R'', s'', t'', c'')$ を、$(U', R', s', t', c')$ と点
$u \in U'$ から始めた対応する問題の解とする。合成
$$
g'' = g \circ g' : U'' \longrightarrow U
$$
が元の問題の解であると主張する。(1), (2), (3), (5), (6) は直ちに従う。(4) については、
射
$$
h'' = s \circ \text{pr}_1 : U'' \times_{g'', U, t} R \longrightarrow U
$$
が補題 \ref{more-groupoids-lemma-double-restrict-property} により局所有限表示かつ Cohen--Macaulay である
（Cohen--Macaulay 射が標的上 fppf 局所的であることには『射の続論』補題
\ref{more-morphisms-lemma-CM-local-source-and-target} を用いる）。
\end{proof}

\noindent
安定化群スキームのファイバーが次元 0 である場合、次のスライス補題を得る。

\begin{lemma}
\label{more-groupoids-lemma-max-slice-quasi-finite}
$S$ をスキームとし、$(U, R, s, t, c, e, i)$ を $S$ 上の亜群スキームとする。
$G \to U$ を安定化群スキームとする。$s$ と $t$ は Cohen--Macaulay かつ局所有限表示で
あると仮定する。$u \in U$ をスキーム $U$ の有限型点とする。『射』定義
\ref{morphisms-definition-finite-type-point} を参照せよ。$G \to U$ は局所準有限であると
仮定する。状況 \ref{more-groupoids-situation-slice} の記法で、次を満たすアフィンスキーム $U'$ と射
$g : U' \to U$ が存在する。
\begin{enumerate}
\item $g$ は埋め込みである。
\item $u \in U'$ である。
\item $g$ は局所有限表示である。
\item 射 $h : U' \times_{g, U, t} R \longrightarrow U$ は平坦、局所有限表示、かつ
局所準有限である。
\item 射 $s', t' : R' \to U'$ は平坦、局所有限表示、かつ局所準有限である。
\end{enumerate}
\end{lemma}

\begin{proof}
補題 \ref{more-groupoids-lemma-max-slice} のとおりに $g : U' \to U$ を取る。
$h^{-1}(u) = F'_u$ なので、$h$ は $(u, e(u))$ で相対次元 $\leq 0$ である。
したがって注意 \ref{more-groupoids-remark-local-source-apply} により、$u \in U''$ を満たす開部分スキーム
$U'' \subset U'$ であって、$U'' \times_{g, U, t} R$ が
$U' \times_{g, U, t} R$ のうち $h$ の相対次元が $\leq 0$ となる最大の開部分スキームで
あるものを得る。$U'$ を $U''$ で置き換えると、$h$ の相対次元は $\leq 0$ となる。
『射』補題 \ref{morphisms-lemma-locally-quasi-finite-rel-dimension-0} により、これは $h$ が
局所準有限であることを意味する。$h$ はなお局所有限表示かつ Cohen--Macaulay なので、
平坦、局所有限表示、かつ局所準有限である。すなわち (4) が成り立つ。
補題 \ref{more-groupoids-lemma-restrict-property} により、その基底変換である $s'$ も平坦、局所有限表示、
かつ局所準有限である。
\end{proof}







\section{亜群のエタール（\'Etale）局所化}
\label{more-groupoids-section-etale-localize}

\noindent
本節では、『射の続論』節 \ref{more-morphisms-section-etale-localization} のエタール（\'etale）局所化の
手法を亜群スキームに適用し始める。この種のより進んだ内容は『空間における亜群の続論』節
\ref{spaces-more-groupoids-section-etale-localize} にある。補題
\ref{more-groupoids-lemma-quasi-finite-over-base-j-proper} は、スキーム上分離的かつ準有限な代数空間に関する
結果、すなわち『空間の射』命題
\ref{spaces-morphisms-proposition-locally-quasi-finite-separated-over-scheme} とその系である
『空間の射』補題 \ref{spaces-morphisms-lemma-locally-quasi-finite-separated-representable} の証明に用いる。

\begin{lemma}
\label{more-groupoids-lemma-quasi-finite-over-base}
$S$ をスキームとし、$(U, R, s, t, c)$ を $S$ 上の亜群スキームとする。
$p \in S$ を点とし、$u \in U$ を $p$ の上にある点とする。次を仮定する。
\begin{enumerate}
\item $U \to S$ は局所有限型である。
\item $U \to S$ は $u$ で準有限である。
\item $U \to S$ は分離的である。
\item $R \to S$ は分離的である。
\item $s$, $t$ は平坦かつ局所有限表示である。
\item $s^{-1}(\{u\})$ は有限である。
\end{enumerate}
このとき、$\kappa(p) = \kappa(p')$ を満たすエタール（\'etale）近傍
$(S', p') \to (S, p)$ と基底変換図式
$$
\xymatrix{
R' \amalg W'
\ar@{=}[r] &
S' \times_S R
\ar[r] \ar@<2ex>[d]^{s'} \ar@<-2ex>[d]_{t'} &
R \ar@<1ex>[d]^s \ar@<-1ex>[d]_t \\
U' \amalg W
\ar@{=}[r] &
S' \times_S U
\ar[r] \ar[d] &
U \ar[d] \\
 &
S' \ar[r] &
S
}
$$
が存在する。ここで等号は開かつ閉な部分スキームへの分解であり、次を満たす。
\begin{enumerate}
\item[(a)] $U'$ には $U$ の $u$ に写る点 $u'$ が存在する。
\item[(b)] 集合論的にファイバー $(U')_{p'}$ は
$t'\big((s')^{-1}(\{u'\})\big)$ に等しい。
\item[(c)] 集合論的にファイバー $(R')_{p'}$ は
$(s')^{-1}\big((U')_{p'}\big)$ に等しい。
\item[(d)] スキーム $U'$ と $R'$ は $S'$ 上有限である。
\item[(e)] $s'(R') \subset U'$ および $t'(R') \subset U'$ である。
\item[(f)] $c'$ を $c$ の基底変換とすると、
$c'(R' \times_{s', U', t'} R') \subset R'$ である。
\item[(g)] 射 $s', t', c'$ は、系
$(U', R', s'|_{R'}, t'|_{R'}, c'|_{R' \times_{s', U', t'} R'})$
を取ることにより亜群構造を定める。
\end{enumerate}
\end{lemma}

\begin{proof}
$f : U \to S$ で $U$ の構造射を表す。仮定 (6) により
$s^{-1}(\{u\}) = \{r_1, \ldots, r_n\}$ と書ける。この集合は有限なので、
$s$ はこれら有限個の逆像の各点で準有限である。『射』補題
\ref{morphisms-lemma-finite-fibre} を参照せよ。したがって
$f \circ s : R \to S$ は各 $r_i$ で準有限である（『射』補題
\ref{morphisms-lemma-composition-quasi-finite}）。ゆえに $r_i$ はファイバー $R_p$ の孤立点である
（『射』補題 \ref{morphisms-lemma-quasi-finite-at-point-characterize}）。
$t(\{r_1, \ldots, r_n\}) = \{u_1, \ldots, u_m\}$ と書く。$m < n$ となることもあり、
$u \in \{u_1, \ldots, u_m\}$ であることに注意する。$t$ は平坦かつ局所有限表示なので、
ファイバーの射 $t_p : R_p \to U_p$ は平坦かつ局所有限表示であり（『射』補題
\ref{morphisms-lemma-base-change-flat} および
\ref{morphisms-lemma-base-change-finite-presentation}）、したがって開である（『射』補題
\ref{morphisms-lemma-fppf-open}）。各 $r_i$ が $R_p$ の孤立点であることから、各
$u_j = t(r_i)$ は $U_p$ の孤立点である。再び『射』補題
\ref{morphisms-lemma-quasi-finite-at-point-characterize} を用いると、$f$ は
$u_1, \ldots, u_m$ で準有限である。

\medskip\noindent
$F_u = s^{-1}(u)$ および $F_{u_j} = s^{-1}(u_j)$ でスキーム論的ファイバーを表す。
$F_u$ は $\kappa(u)$ 上局所有限型で点が有限個なので、$\kappa(u)$ 上有限である
（例えば、より一般的な『射』補題
\ref{morphisms-lemma-locally-quasi-finite-qc-source-universally-bounded} から従う）。
補題 \ref{more-groupoids-lemma-two-fibres} により、$F_u$ と $F_{u_j}$ は $\kappa(u)$ と
$\kappa(u_j)$ の共通の体拡大上で同型になる。したがって $F_{u_j}$ は
$\kappa(u_j)$ 上有限である。特に各 $j = 1, \ldots, m$ に対して
$s^{-1}(\{u_j\})$ は有限集合である。したがって各 $u_j$ に対しても仮定 (2), (6) が
成り立つ（$U \to S$ が $u_j$ で準有限であることは上で示した）。そこで第1段落の議論を
各 $u_j$ に適用すると、$R \to U$ は集合
$$
\{r_1, \ldots, r_N\} = s^{-1}(\{u_1, \ldots, u_m\})
$$
の各点で準有限である。また
$t(\{r_1, \ldots, r_N\}) = \{u_1, \ldots, u_m\}$ および
$t^{-1}(\{u_1, \ldots, u_m\}) = \{r_1, \ldots, r_N\}$ である。これは $R$ が
亜群を成すからである\footnote{亜群の言葉で説明すると、元の集合
$\{r_1, \ldots, r_n\}$ は始点が $u$ である射の集合であった。
$\{u_1, \ldots, u_m\}$ は $u$ と同型な対象の集合であり、
$\{r_1, \ldots, r_N\}$ は $u$ と同値なすべての対象の間の全射の集合である。}。さらに
$\text{pr}_0(c^{-1}(\{r_1, \ldots, r_N\})) = \{r_1, \ldots, r_N\}$
および
$\text{pr}_1(c^{-1}(\{r_1, \ldots, r_N\})) = \{r_1, \ldots, r_N\}$
である。同様に
$e(\{u_1, \ldots, u_m\}) \subset \{r_1, \ldots, r_N\}$ および
$i(\{r_1, \ldots, r_N\}) = \{r_1, \ldots, r_N\}$ を得る。

\medskip\noindent
『射の続論』補題 \ref{more-morphisms-lemma-etale-splits-off-quasi-finite-part-technical} を
組 $(U \to S, \{u_1, \ldots, u_m\})$ および
$(R \to S, \{r_1, \ldots, r_N\})$ に適用すると、エタール（\'etale）近傍
$(S', p') \to (S, p)$ を得る。
これは同一視 $\kappa(p) = \kappa(p')$ を誘導し、$S' \times_S U$ と
$S' \times_S R$ は
$$
S' \times_S U = U' \amalg W, \quad
S' \times_S R = R' \amalg W'
$$
と分解する。ここで $U' \to S'$ は有限で、$(U')_{p'}$ は
$\{u_1, \ldots, u_m\}$ に全単射で写り、$R' \to S'$ は有限で、$(R')_{p'}$ は
$\{r_1, \ldots, r_N\}$ に全単射で写る。さらに $W_{p'}$（resp.\ $(W')_{p'}$）の点は
どの $u_j$（resp.\ $r_i$）にも写らない。この時点で補題の (a), (b), (c), (d) が成り立つ。
また (e), (f) の包含は $p'$ 上のファイバーで成り立つ。すなわち
$s'((R')_{p'}) \subset (U')_{p'}$、$t'((R')_{p'}) \subset (U')_{p'}$、および
$c'((R' \times_{s', U', t'} R')_{p'}) \subset (R')_{p'}$ である。

\medskip\noindent
$S'$ を $p'$ の Zariski 開近傍で置き換え、(e), (f) の包含を成り立たせられると主張する。
例えば集合 $E = (s'|_{R'})^{-1}(W)$ を考える。これは $R'$ の開かつ閉な部分集合で、
$p'$ の上にある $R'$ の点を含まない。$R' \to S'$ は閉なので、$S'$ を
$S' \setminus (R' \to S')(E)$ で置き換えると $E$ が空である状況になる。言い換えれば、
$s'$ は $R'$ を $U'$ に写す。この性質は $S'$ をさらに縮小しても保たれる。同様に、$t'$ が
$R'$ を $U'$ に写るようにできる。これで (e) が成り立つ。同様に集合
$E = (c'|_{R' \times_{s', U', t'} R'})^{-1}(W')$ を考える。これは $S'$ 上有限なスキーム
$R' \times_{s', U', t'} R'$ の開かつ閉な部分集合で、$p'$ の上の点を含まない。
したがって $S'$ を
$S' \setminus (R' \times_{s', U', t'} R' \to S')(E)$
で置き換えると $E$ は空となり、(f) の包含を得る。恒等射
$e' : S' \times_S U \to S' \times_S R$ と逆射
$i' : S' \times_S R \to S' \times_S R$ にも同じ議論を繰り返し、$S'$ をさらに縮小した後
$(e'|_{U'})^{-1}(W') = \emptyset$ および $(i'|_{R'})^{-1}(W') = \emptyset$ と仮定できる。

\medskip\noindent
この時点で構造
$$
(U', R', s'|_{R'}, t'|_{R'}, c'|_{R' \times_{t', U', s'} R'},
e'|_{U'}, i'|_{R'}).
$$
を考えられる。$S'$ 上の亜群スキームの公理は、亜群スキーム
$(S' \times_S U, S' \times_S R, s', t', c', e', i')$ に対して成り立つので、ここでも成り立つ。
\end{proof}

\begin{lemma}
\label{more-groupoids-lemma-quasi-finite-over-base-j-proper}
$S$ をスキームとし、$(U, R, s, t, c)$ を $S$ 上の亜群スキームとする。
$p \in S$ を点とし、$u \in U$ を $p$ の上にある点とする。補題
\ref{more-groupoids-lemma-quasi-finite-over-base} の仮定 (1) -- (6) に加えて次を仮定する。
\begin{enumerate}
\item[(7)] $j : R \to U \times_S U$ は普遍的閉である\footnote{他の条件を考慮すると、
これは $j$ が固有であることを要求するのと同値である。}。
\end{enumerate}
このとき、補題 \ref{more-groupoids-lemma-quasi-finite-over-base} の (a) -- (g) と次を満たす
$(S', p') \to (S, p)$、分解 $S' \times_S U = U' \amalg W$ と
$S' \times_S R = R' \amalg W'$、および $u' \in U'$ を選べる。
\begin{enumerate}
\item[(h)] $R'$ は $U'$ への $S' \times_S R$ の制限である。
\end{enumerate}
\end{lemma}

\begin{proof}
スキーム $S$ 上の亜群 $(U, R, s, t, c)$ と点 $p$, $u$ に補題
\ref{more-groupoids-lemma-quasi-finite-over-base} を適用する。これによりエタール（\'etale）近傍
$(S', p') \to (S, p)$、非交和分解
$$
S' \times_S U = U' \amalg W, \quad
S' \times_S R = R' \amalg W'
$$
および結論 (a), (b), (c), (d), (e), (f), (g) を満たす $u' \in U'$ を得る。
結論 (a) -- (g) に影響を与えずに、$S'$ を $p'$ のより小さい近傍へ縮小できる。
適切に縮小すれば結論 (h) も成り立つことを示す。$j'$ で $j$ の $S'$ への基底変換を表す。
結論 (e) より、ある開かつ閉な $Rest$ 部分に対して
$$
j'^{-1}(U' \times_{S'} U') = R' \amalg Rest
$$
である。結論 (d) により $U' \to S'$ は有限なので、$U' \times_{S'} U'$ は $S'$ 上有限である。
$j$ は普遍的閉なので $j'$ も普遍的閉であり、したがって $j'|_{Rest}$ も普遍的閉である。
結論 (b), (c) により
$$
(U' \times_{S'} U' \to S') \circ j'|_{Rest} :
Rest
\longrightarrow
S'
$$
の $p'$ 上のファイバーは空である。$Rest \to S'$ は閉射の合成として閉だから、$S'$ を
$S' \setminus \Im(Rest \to S')$ で置き換え、$Rest = \emptyset$ と仮定できる。
これはまさに、$R'$ が開部分スキーム $U' \subset S' \times_S U$ への
$S' \times_S R$ の制限であるという条件である。『亜群スキーム』補題
\ref{groupoids-lemma-restrict-groupoid-relation} とその証明を参照せよ。
\end{proof}






\section{有限亜群}
\label{more-groupoids-section-finite-groupoids}

\noindent
亜群スキーム $(U, R, s, t, c)$ は、射 $s$, $t$ が有限であるとき、しばしば
{\it 有限}と呼ばれる。しかしこれは $U$、$R$、あるいは商層 $U/R$ が何かの上で有限で
あることを意味しないため、混同を招く可能性がある。

\begin{lemma}
\label{more-groupoids-lemma-finite-stratify}
$(U, R, s, t, c)$ をスキーム $S$ 上の亜群スキームとし、$s, t$ は有限であると仮定する。
$R$-不変な閉部分スキームの列
$$
U = Z_0 \supset Z_1 \supset Z_2 \supset \ldots
$$
であって、$\bigcap Z_r = \emptyset$ であり、
$s^{-1}(Z_{r - 1}) \setminus s^{-1}(Z_r) \to Z_{r - 1} \setminus Z_r$
が階数 $r$ の有限局所自由射となるものが存在する。
\end{lemma}

\begin{proof}
$\{Z_r\}$ を、有限型準連接加群 $s_*\mathcal{O}_R$ の Fitting イデアルが与える $U$ の
層別化とする。『因子』補題 \ref{divisors-lemma-locally-free-rank-r-pullback} を参照せよ。
恒等射 $e : U \to R$ は $s$ の切断なので、$s_*\mathcal{O}_R$ は $\mathcal{O}_S$ を
直和因子として含む。したがって $U = Z_{-1} = Z_0$ である（詳細は省略する）。
Fitting イデアルの形成は基底変換と可換なので（『代数の続論』補題
\ref{more-algebra-lemma-fitting-ideal-basics}）、図式 (\ref{more-groupoids-equation-diagram}) の右下の正方形が
Cartesian であることから、$s^{-1}(Z_r)$ は
$\text{pr}_{1, *}\mathcal{O}_{R \times_{s, U, t} R}$ の第 $r$ Fitting イデアルに対応する。
左下の正方形も Cartesian であることを用いると $s^{-1}(Z_r) = t^{-1}(Z_r)$ を得る。
言い換えれば $Z_r$ は $R$-不変である。射
$s^{-1}(Z_{r - 1}) \setminus s^{-1}(Z_r) \to Z_{r - 1} \setminus Z_r$
は階数 $r$ の有限局所自由射である。実際、『因子』補題
\ref{divisors-lemma-locally-free-rank-r-pullback} により、加群 $s_*\mathcal{O}_R$ は
$Z_{r - 1} \setminus Z_r$ 上で階数 $r$ の有限局所自由加群に引き戻される。
\end{proof}

\begin{lemma}
\label{more-groupoids-lemma-finite-flat-over-almost-dense-subscheme}
$(U, R, s, t, c)$ をスキーム $S$ 上の亜群スキームとし、$s, t$ は有限であると仮定する。
開部分スキーム $W \subset U$ と閉部分スキーム $W' \subset W$ であって次を満たすものが存在する。
\begin{enumerate}
\item $W$ と $W'$ は $R$-不変である。
\item 集合論的に $U = t(s^{-1}(\overline{W}))$ である。
\item $W$ は $W'$ の厚化である。
\item 制限 $(W', R', s', t', c')$ の写像 $s'$, $t'$ は有限局所自由である。
\end{enumerate}
\end{lemma}

\begin{proof}
補題 \ref{more-groupoids-lemma-finite-stratify} の層別化
$U = Z_0 \supset Z_1 \supset Z_2 \supset \ldots$ を考える。

\medskip\noindent
非交和 $W = \coprod_{r \geq 1} W_r$ と $W' = \coprod_{r \geq 1} W'_r$ を構成する。
ここで各 $W'_r \to W_r$ は $U$ の $R$-不変な部分スキームの厚化であり、制限
$(W_r', R_r', s_r', t_r', c_r')$ の射 $s_r', t_r'$ は階数 $r$ の有限局所自由射とする。
まず $W_1 = W'_1 = U \setminus Z_1$ とおく。これは $U$ の $R$-不変な開部分スキームで、
$W_0$ は $W'_0$ の厚化であり、制限 $(W_1', R_1', s_1', t_1', c_1')$ の写像
$s_1'$, $t_1'$ は同型、すなわち階数 $1$ の有限局所自由射である。さらに
$U \setminus Z_1$ の各点は $t(s^{-1}(\overline{W_1}))$ に含まれる。

\medskip\noindent
$r \leq n$ に対して部分スキーム $W'_r \subset W_r \subset U$ が構成され、次を満たすと仮定する。
\begin{enumerate}
\item $W_1, \ldots, W_n$ は互いに素である。
\item $W_r$ と $W_r'$ は $R$-不変である。
\item 集合論的に
$U \setminus Z_n \subset \bigcup_{r \leq n} t(s^{-1}(\overline{W_r}))$ である。
\item $W_r$ は $W'_r$ の厚化である。
\item 制限 $(W_r', R_r', s_r', t_r', c_r')$ の写像 $s_r'$, $t_r'$ は階数 $r$ の
有限局所自由射である。
\end{enumerate}
このとき集合論的に
$$
W_{n + 1} = Z_n \setminus
\left(
Z_{n + 1} \cup \bigcup\nolimits_{r \leq n} t(s^{-1}(\overline{W_r}))
\right)
$$
とおき、スキーム論的に
$$
W'_{n + 1} = Z_n \setminus
\left(
Z_{n + 1} \cup \bigcup\nolimits_{r \leq n} t(s^{-1}(\overline{W_r}))
\right)
$$
とおく。このとき $W_{n + 1}$ は $U$ の $R$-不変な開部分スキームである。実際、
$Z_{n + 1} \setminus \overline{U \setminus Z_{n + 1}}$ は $U$ で開であり、条件 (3) と $t$ が
閉射であることから、$\overline{U \setminus Z_{n + 1}}$ は取り除く閉集合
$\bigcup\nolimits_{r \leq n} t(s^{-1}(\overline{W_r}))$ に含まれる。
$W'_{n + 1}$ は $W_{n + 1}$ の閉部分スキームで、同じ台位相空間をもつことは明らかである。
最後に、性質 (1), (2), (3) は明らかであり、性質 (5) は補題
\ref{more-groupoids-lemma-finite-stratify} から従う。

\medskip\noindent
補題 \ref{more-groupoids-lemma-finite-stratify} により $\bigcap Z_r = \emptyset$ である。したがって $U$ の
各点は、ある $n$ に対して $U \setminus Z_n$ に含まれる。ゆえに集合論的に
$U = \bigcup_{r \geq 1} t(s^{-1}(\overline{W_r}))$ であり、(2) が成り立つ。
以上より $W' \subset W$ は (1), (2), (3), (4) を満たす。
\end{proof}

\noindent
$(U, R, s, t, c)$ を亜群スキームとする。点 $u \in U$ に対し、$u$ の
{\it $R$-軌道}とは $U$ の部分集合 $t(s^{-1}(\{u\}))$ のことである。

\begin{lemma}
\label{more-groupoids-lemma-finite-flat-over-almost-dense-subscheme-addendum}
補題 \ref{more-groupoids-lemma-finite-flat-over-almost-dense-subscheme} において、さらに $s$, $t$ が
有限表示であると仮定する。このとき
\begin{enumerate}
\item 射 $W' \to W$ は有限表示である。
\item $u \in U$ の $R$-軌道が $U$ の既約成分の生成点からなるならば $u \in W$ である。
\end{enumerate}
\end{lemma}

\begin{proof}
この場合、補題 \ref{more-groupoids-lemma-finite-stratify} の層別化
$U = Z_0 \supset Z_1 \supset Z_2 \supset \ldots$ は有限表示の閉埋め込み $Z_k \to U$ に
よって与えられる。『因子』補題 \ref{divisors-lemma-locally-free-rank-r-pullback} を参照せよ。
$W' \to W$ は局所的に開集合 $W$ と $Z_r$ の交わりで与えられるので、(1) は直ちに従う。
(2) を示すため、$u$ の軌道を $\{u_1, \ldots, u_n\}$ とする。閉部分スキーム $Z_k$ は
$R$-不変で $\bigcap Z_k = \emptyset$ だから、すべての $i$ に対して $u_i \in Z_k$ かつ
$u_i \not \in Z_{k + 1}$ となる $k$ がある。$Z_k \to U$ と $Z_{k + 1} \to U$ の像は
局所構成可能である（『射』定理 \ref{morphisms-theorem-chevalley}）。$u_i \in U$ は $U$ の
既約成分の生成点なので、集合論的に $Z_k \setminus Z_{k + 1}$ に含まれる $u_i$ の開近傍
$U_i$ が存在する（『性質』補題 \ref{properties-lemma-generic-point-in-constructible}）。
補題 \ref{more-groupoids-lemma-finite-flat-over-almost-dense-subscheme} の証明では、$W$ を
$W_r \subset Z_{r - 1} \setminus Z_r$ かつ
$U = \bigcup t(s^{-1}(\overline{W_r}))$ を満たす非交和 $\coprod W_r$ として構成した。
$\{u_1, \ldots, u_n\}$ は $R$-軌道なので、$u \in t(s^{-1}(\overline{W_r}))$ ならば
ある $i$ に対して $u_i \in \overline{W_r}$ であり、したがって
$U_i \cap W_r \not = \emptyset$、ゆえに $r = k$ である。以上より、所望のとおり $u$ は
$$
W_{k + 1} = Z_k \setminus
\left(
Z_{k + 1} \cup \bigcup\nolimits_{r \leq k} t(s^{-1}(\overline{W_r}))
\right)
$$
に含まれる。
\end{proof}

\begin{lemma}
\label{more-groupoids-lemma-invariant-affine-open-around-generic-point}
$(U, R, s, t, c)$ をスキーム $S$ 上の亜群スキームとする。$s, t$ は有限かつ有限表示で、
$U$ は準分離であると仮定する。$u_1, \ldots, u_m \in U$ を、その軌道が $U$ の既約成分の
生成点からなる点とする。このとき、$R$-不変な部分スキーム $V' \subset V \subset U$ で
次を満たすものが存在する。
\begin{enumerate}
\item $u_1, \ldots, u_m \in V'$ である。
\item $V$ は $U$ で開である。
\item $V'$ と $V$ はアフィンである。
\item $V' \subset V$ は有限表示の厚化である。
\item 制限 $(V', R', s', t', c')$ の射 $s', t'$ は有限局所自由である。
\end{enumerate}
\end{lemma}

\begin{proof}
$W' \subset W \subset U$ を補題 \ref{more-groupoids-lemma-finite-flat-over-almost-dense-subscheme} のとおりに取る。
補題 \ref{more-groupoids-lemma-finite-flat-over-almost-dense-subscheme-addendum} により $u_j \in W$ であり、
$W' \to W$ は有限表示の厚化である。『極限』補題
\ref{limits-lemma-affines-glued-in-closed-affine} により、$u_j$ を含む $W'$ の $R$-不変な
アフィン開部分スキーム $V'$ を見つければ十分である（そのとき対応する開部分スキーム
$V \subset W$ はアフィンとなる）。そこで $(U, R, s, t, c)$ を $W'$ への制限
$(W', R', s', t', c')$ で置き換えてよい。言い換えれば、$s$, $t$ が有限局所自由である
亜群スキーム $(U, R, s, t, c)$ を考えていると仮定できる。『性質』補題
\ref{properties-lemma-maximal-points-affine} により $u_1, \ldots, u_m$ の軌道の合併を含む
アフィン開集合が存在する。最後に『亜群スキーム』補題
\ref{groupoids-lemma-find-invariant-affine} を適用して結論を得る。
\end{proof}

\noindent
次の補題は補題 \ref{more-groupoids-lemma-invariant-affine-open-around-generic-point} の特別な場合だが、
この場合は議論が少し容易なので改めて証明する（補題
\ref{more-groupoids-lemma-finite-flat-over-almost-dense-subscheme-addendum} を使わずに済む）。

\begin{lemma}
\label{more-groupoids-lemma-invariant-affine-open-around-generic-point-Noetherian}
$(U, R, s, t, c)$ をスキーム $S$ 上の亜群スキームとする。$s, t$ は有限で、$U$ は
局所 Noether であると仮定し、$u_1, \ldots, u_m \in U$ を、その軌道が $U$ の既約成分の
生成点からなる点とする。このとき、次を満たす $R$-不変な部分スキーム
$V' \subset V \subset U$ が存在する。
\begin{enumerate}
\item $u_1, \ldots, u_m \in V'$ である。
\item $V$ は $U$ で開である。
\item $V'$ と $V$ はアフィンである。
\item $V' \subset V$ は厚化である。
\item 制限 $(V', R', s', t', c')$ の射 $s', t'$ は有限局所自由である。
\end{enumerate}
\end{lemma}

\begin{proof}
$\{u_{j1}, \ldots, u_{jn_j}\}$ を $u_j$ の軌道とする。
$W' \subset W \subset U$ を補題 \ref{more-groupoids-lemma-finite-flat-over-almost-dense-subscheme} のとおりに取る。
$U = t(s^{-1}(\overline{W}))$ なので、少なくとも一つの $u_{ji} \in \overline{W}$ がある。
$u_{ji}$ は既約成分の生成点で $U$ は局所 Noether だから、$u_{ji} \in W$ である。
$W$ は $R$-不変なので $u_j \in W$ であり、実際に軌道全体が $W$ に含まれる。
『スキームのコホモロジー』補題
\ref{coherent-lemma-image-affine-finite-morphism-affine-Noetherian} により、
$u_1, \ldots, u_m$ を含む $W'$ の $R$-不変なアフィン開部分スキーム $V'$ を見つければ
十分である（そのとき対応する開部分スキーム $V \subset W$ はアフィンとなる）。そこで
$(U, R, s, t, c)$ を $W'$ への制限 $(W', R', s', t', c')$ で置き換えてよい。
言い換えれば、$s$, $t$ が有限局所自由である亜群スキーム $(U, R, s, t, c)$ を考えていると
仮定できる。『性質』補題 \ref{properties-lemma-maximal-points-affine} により
$\{u_{ij}\}$ を含むアフィン開集合が存在する（局所 Noether スキームは準分離である。
『性質』補題 \ref{properties-lemma-locally-Noetherian-quasi-separated} を参照せよ）。
最後に『亜群スキーム』補題 \ref{groupoids-lemma-find-invariant-affine} を適用して結論を得る。
\end{proof}

\begin{lemma}
\label{more-groupoids-lemma-find-affine-integral}
$(U, R, s, t, c)$ をスキーム $S$ 上の亜群スキームとし、$s, t$ は整であるとする。
$g : U' \to U$ を整射で、$U$ の各 $R$-軌道が $g(U')$ と交わるものとする。
$(U', R', s', t', c')$ を $U'$ への $R$ の制限とする。$u' \in U'$ が $R'$-不変な
アフィン開集合に含まれるならば、その像 $u \in U$ は $U$ の $R$-不変なアフィン開集合に
含まれる。
\end{lemma}

\begin{proof}
$W' \subset U'$ を $R'$-不変なアフィン開集合とする。
$\tilde R = U' \times_{g, U, t} R$ とおき、写像
$\text{pr}_0 : \tilde R \to U'$ および $h = s \circ \text{pr}_1 : \tilde R \to U$ を考える。
$\text{pr}_0$ と $h$ は整である。したがって $\tilde W = \text{pr}_0^{-1}(W')$ はアフィンである。
$W'$ は $R'$-不変なので、像 $W = h(\tilde W)$ は集合論的に $R$-不変で、集合論的に
$\tilde W = h^{-1}(W)$ である（詳細は省略する）。そこで $W$ が開であることを示せれば、
$W$ はスキームで、射 $\tilde W \to W$ は整かつ全射となり、『極限』命題
\ref{limits-proposition-affine} により $W$ はアフィンである。一方、各軌道が $U'$ と交わるという
仮定から $h : \tilde R \to U$ は全射である。整全射は商位相を与えるので（『位相』補題
\ref{topology-lemma-closed-morphism-quotient-topology} および『射』補題
\ref{morphisms-lemma-integral-universally-closed}）、$W$ は開である。
\end{proof}

\noindent
次の技術的補題は、準アフィンスキーム上の有限亜群という状況で「ほぼ」不変な関数を構成する。

\begin{lemma}
\label{more-groupoids-lemma-find-almost-invariant-function}
$(U, R, s, t, c)$ を、$s, t$ が有限かつ有限表示である亜群スキームとする。
$u_1, \ldots, u_m \in U$ を、その $R$-軌道が $U$ の既約成分の生成点からなる点とする。
$j : U \to \Spec(A)$ を埋め込みとする。$I \subset A$ を
$j(U) \cap V(I) = \emptyset$ かつ $V(I) \cup j(U)$ が $\Spec(A)$ で閉となるイデアルとする。
このとき $h \in I$ であって、$j^{-1}D(h)$ が $u_1, \ldots, u_m$ を含む $U$ の
$R$-不変なアフィン開部分スキームとなるものが存在する。
\end{lemma}

\begin{proof}
$u_1, \ldots, u_m \in V' \subset V \subset U$ を補題
\ref{more-groupoids-lemma-invariant-affine-open-around-generic-point} のとおりに取る。
$U \setminus V$ は $U$ で閉、$j$ は埋め込み、$V(I) \cup j(U)$ は $\Spec(A)$ で閉なので、
$V(J) = V(I) \cup j(U \setminus V)$ を満たすイデアル $J \subset I$ が存在する。
例えば $j(U \setminus V)$ 上で消える $I$ の元からなるイデアルを取ればよい。したがって
$(U, R, s, t, c)$、$j : U \to \Spec(A)$、$I$ をそれぞれ
$(V', R', s', t', c')$、$j|_{V'} : V' \to \Spec(A)$、$J$ で置き換えられる。
言い換えれば、$U$ はアフィンで、$s$, $t$ は有限局所自由であると仮定してよい。
$u_1, \ldots, u_m$ の $R$-軌道のすべての点で消えない任意の $f \in I$ を取る
（『代数』補題 \ref{algebra-lemma-silly}）。
$$
g = \text{Norm}_s(t^\sharp(j^\sharp(f))) \in \Gamma(U, \mathcal{O}_U)
$$
を考える。$f \in I$ で $V(I) \cup j(U)$ は閉なので、$U \cap D(f) \to D(f)$ は閉埋め込みで
ある。したがって、ある $n > 0$ に対して $f^ng$ はある元 $h \in I$ の像である。この $h$ が
所望のものであると主張する。実際、『亜群スキーム』補題
\ref{groupoids-lemma-determinant-trick} で $g$ は $R$-不変関数であることを示したので、
$D(g) \subset U$ は $R$-不変である。$f$ は $u_j$ の軌道上で消えないから、関数 $g$ は
$u_j$ で消えない。さらに $V(g) \supset V(j^\sharp(f))$ であり、したがって
$j^{-1}D(h) = D(g)$ である。
\end{proof}

\begin{lemma}
\label{more-groupoids-lemma-no-specializations-map-to-same-point}
$(U, R, s, t, c)$ を亜群スキームとする。$s, t$ が有限で、$u, u' \in R$ が同じ軌道にある
相異なる点ならば、$u'$ は $u$ の特殊化ではない。
\end{lemma}

\begin{proof}
$s(r) = u$ および $t(r) = u'$ を満たす $r \in R$ を取る。$u \leadsto u'$ なら、
$s(r') = u'$ を満たす非自明な特殊化 $r \leadsto r'$ が存在する。『スキーム』補題
\ref{schemes-lemma-quasi-compact-closed} を参照せよ。$u'' = t(r')$ とおく。有限射の
ファイバー内には特殊化がないので $u'' \not = u'$ である。したがって続けて、
$s(r'') = u''$ を満たす非自明な特殊化 $r' \leadsto r''$ などを見つけられる。
これは $u$ の軌道が特殊化の無限列
$u \leadsto u' \leadsto u'' \leadsto \ldots$
を含むことを示すが、軌道 $t(s^{-1}(\{u\}))$ は有限なので矛盾である。
\end{proof}

\begin{lemma}
\label{more-groupoids-lemma-get-affine}
$j : V \to \Spec(A)$ をスキームの準コンパクトな埋め込みとする。
$f \in A$ は $j^{-1}D(f)$ がアフィンで、$j(V) \cap V(f)$ が閉となるものとする。
このとき $V$ はアフィンである。
\end{lemma}

\begin{proof}
これは『射』補題 \ref{morphisms-lemma-get-affine} から従うが、直接の証明も与える。
$A' = \Gamma(V, \mathcal{O}_V)$ とおく。このとき
$j' : V \to \Spec(A')$ は準コンパクトな開埋め込みである。『性質』補題
\ref{properties-lemma-quasi-affine} を参照せよ。$f' \in A'$ を $f$ の像とする。
このとき $(j')^{-1}D(f') = j^{-1}D(f)$ はアフィンである。一方、$j'(V) \cap V(f')$ は
$\Spec(A')$ の部分スキームで、$\Spec(A)$ の閉部分スキーム $j(V) \cap V(f)$ に同型に写る。
したがって、例えば『スキーム』補題 \ref{schemes-lemma-section-immersion} により、これは
$\Spec(A')$ で閉である。そこで $A$ を $A'$ で置き換え、$j$ は開埋め込みで
$A = \Gamma(V, \mathcal{O}_V)$ であると仮定してよい。

\medskip\noindent
この場合 $j(V) = \Spec(A)$ であると主張する。これで証明は終わる。そうでなければ、補集合と
交わり、閉部分集合 $j(V) \cap V(f)$ を避ける主アフィン開集合
$D(g) \subset \Spec(A)$ が存在する。『性質』補題
\ref{properties-lemma-invert-f-affine} により、$j$ は $j^{-1}D(f)$ を $D(f)$ に同型に写す。
したがって $D(g)$ は $V(f)$ と交わる。一方、$j^{-1}D(g)$ はアフィン開集合
$j^{-1}D(f)$ の主開集合なのでアフィンである。再び『性質』補題
\ref{properties-lemma-invert-f-affine} により、$D(g)$ は
$j^{-1}D(g) \subset j^{-1}D(f)$ と同型である。これは $D(g) \subset D(f)$ を意味し、矛盾である。
\end{proof}

\begin{lemma}
\label{more-groupoids-lemma-find-affine-codimension-1}
$(U, R, s, t, c)$ を亜群スキームとし、$u \in U$ とする。次を仮定する。
\begin{enumerate}
\item $s, t$ は有限射である。
\item $U$ は分離的かつ局所 Noether である。
\item $u$ の軌道にある各点 $u'$ に対して $\dim(\mathcal{O}_{U, u'}) \leq 1$ である。
\end{enumerate}
このとき $u$ は $U$ の $R$-不変なアフィン開集合に含まれる。
\end{lemma}

\begin{proof}
$u$ の $R$-軌道は有限である。条件 (2), (3) により、これは $U$ のアフィン開集合 $U'$ に
含まれる。『多様体』命題
\ref{varieties-proposition-finite-set-of-points-of-codim-1-in-affine} を参照せよ。このとき
$t(s^{-1}(U \setminus U'))$ は $u$ を含まない $U$ の $R$-不変な閉部分集合である。
したがって $U \setminus t(s^{-1}(U \setminus U'))$ は $u$ を含む $U'$ の $R$-不変な
開集合である。$U$ をこの開集合で置き換え、$U$ は準アフィンであると仮定してよい。

\medskip\noindent
補題 \ref{more-groupoids-lemma-find-affine-integral} により $U$ をその被約化で置き換え、$U$ は被約であると
仮定してよい。これは補題 \ref{more-groupoids-lemma-finite-flat-over-almost-dense-subscheme} の
$R$-不変な部分スキーム $W' \subset W \subset U$ が $W' = W$ を満たすことを意味する。
$U = t(s^{-1}(\overline{W}))$ なので、$u$ の $R$-軌道のある点 $u'$ は
$\overline{W}$ に含まれる。補題 \ref{more-groupoids-lemma-find-affine-integral} により $U$ を
$\overline{W}$ で、$u$ を $u'$ で置き換えられる。したがって、稠密開な $R$-不変部分スキーム
$W \subset U$ であって、制限 $(W, R_W, s_W, t_W, c_W)$ の射 $s_W, t_W$ が有限局所自由で
あるものが存在すると仮定してよい。

\medskip\noindent
$u \in W$ なら『亜群スキーム』補題 \ref{groupoids-lemma-find-invariant-affine} により終わる
（$W$ は準アフィンなので、$W$ の任意の有限点集合はアフィン開集合に含まれる。『性質』補題
\ref{properties-lemma-ample-finite-set-in-affine} を参照せよ）。そこで $u \not \in W$ と仮定する。
したがって $u$ の軌道の点はどれも $W$ にない。$\xi \in U$ を、$u$ の軌道にある点 $u'$ への
非自明な特殊化をもつ点とする。$u$ の軌道の点の間には特殊化がないので（補題
\ref{more-groupoids-lemma-no-specializations-map-to-same-point}）、$\xi$ は軌道に含まれない。仮定 (3) により
$\xi$ は $U$ の生成点であり、したがって $\xi \in W$ である。$U$ は Noether なので、
このような点は有限個 $\xi_1, \ldots, \xi_m \in W$ である。$s_W, t_W$ は平坦なので、各
$\xi_j$ の軌道は $W$、したがって $U$ の既約成分の生成点からなる。

\medskip\noindent
$j : U \to \Spec(A)$ を $U$ のアフィンスキームへの埋め込みとする（$U$ は準アフィンなので
可能である）。$J \subset A$ を $V(J) \cap j(W) = \emptyset$ かつ
$V(J) \cup j(W)$ が閉となるイデアルとする。補題
\ref{more-groupoids-lemma-find-almost-invariant-function} を亜群スキーム $(W, R_W, s_W, t_W, c_W)$、射
$j|_W : W \to \Spec(A)$、点 $\xi_j$、イデアル $J$ に適用し、すべての $j$ に対して
$\xi_j$ を含む $R_W$-不変なアフィン開集合となる $(j|_W)^{-1}D(f)$ を与える
$f \in J$ を得る。$f \in J$ なので $j^{-1}D(f) \subset W$ である。すなわち
$j^{-1}D(f)$ はすべての $\xi_j$ を含み $W$ に含まれる、$U$ の $R$-不変なアフィン開集合である。

\medskip\noindent
$Z$ を
$$
U \setminus j^{-1}D(f) = j^{-1}V(f).
$$
上の誘導される被約閉部分スキーム構造とする。このとき $Z$ は集合論的に $R$-不変である
（ただしスキーム論的には $R$-不変でないことがある）。$(Z, R_Z, s_Z, t_Z, c_Z)$ を
$Z$ への $R$ の制限とする。$Z \to U$ は有限なので $s_Z$, $t_Z$ は有限である。
$u \in Z$ なので $u$ の軌道は $Z$ に含まれ、$Z$ の点として見た $u$ の $R_Z$-軌道と
一致する。$\dim(\mathcal{O}_{U, u'}) \leq 1$ で、すべての $j$ に対して
$\xi_j \not \in Z$ なので、$u$ の軌道にあるすべての $u'$ に対して
$\dim(\mathcal{O}_{Z, u'}) \leq 0$ である。言い換えれば、$u$ の $R_Z$-軌道は $Z$ の
既約成分の生成点からなる。

\medskip\noindent
$I \subset A$ を $V(I) \cap j(U) =\emptyset$ かつ $V(I) \cup j(U)$ が閉となるイデアルとする。
補題 \ref{more-groupoids-lemma-find-almost-invariant-function} を亜群スキーム
$(Z, R_Z, s_Z, t_Z, c_Z)$、制限 $j|_Z$、イデアル $I$、点 $u \in Z$ に適用し、
$j^{-1}D(h) \cap Z$ が $u$ を含む $R_Z$-不変なアフィン開集合となる $h \in I$ を得る。

\medskip\noindent
$R_W$-不変な（『亜群スキーム』補題 \ref{groupoids-lemma-determinant-trick}）関数
$$
g = 
\text{Norm}_{s_W}(t_W^\sharp(j^\sharp(h)|_W)) \in \Gamma(W, \mathcal{O}_W)
$$
を考える（以下では $j^{-1}D(f)$ への $g$ の制限しか用いず、この場合ノルムはアフィン間の
有限局所自由射に沿う）。
$$
V = (W_g \cap j^{-1}D(f)) \cup (j^{-1}D(h) \cap Z)
$$
は $U$ の $R$-不変なアフィン開集合であると主張する。これで補題の証明は終わる。構成により
これは集合論的に $R$-不変である。$V$ は構成可能集合なので、開であることを示すには $U$ で
一般化に関して閉じていることを示せば十分である（『位相』補題
\ref{topology-lemma-characterize-closed-Noetherian}、またはより一般に『位相』補題
\ref{topology-lemma-constructible-stable-specialization-closed}）。
$W_g \cap j^{-1}D(f)$ は $U$ で開なので、$u_2 \in j^{-1}D(h) \cap Z$ を満たす $U$ の特殊化
$u_1 \leadsto u_2$ を考えれば十分である。これは $h$ が $j(u_2)$ で非零で、$u_2 \in Z$ で
あることを意味する。$u_1 \in Z$ なら $j(u_1) \leadsto j(u_2)$ であり、$h$ は $j(u_2)$ で
非零なので $j(u_1)$ でも非零である。したがって $u_1 \in V$ である。
$u_1 \not \in Z$ かつ $W_g \cap j^{-1}D(f)$ にも含まれないとする。
$Z = j^{-1}V(f)$ の補集合は $W \cap j^{-1}D(f)$ に含まれるので、$u_1 \in W$ かつ
$u_1 \not \in W_g$ である。したがって $s(r_1) = u_1$ を満たし、$h$ が $t(r_1)$ で
零となる点 $r_1 \in R$ が存在する。$s$ は有限なので、$s(r_2) = u_2$ を満たす特殊化
$r_1 \leadsto r_2$ が存在する。しかしこのとき $h$ は $u'_2 = t(r_2)$ で零となる。これは
$j^{-1}D(h) \cap Z$ が $R$-不変で $u_2$ を含むことに矛盾する。したがって $V$ は開である。

\medskip\noindent
関数 $h \in I$ に対して $V \subset j^{-1}D(h)$ であることに注意する。したがって埋め込み
$$
j' : V \longrightarrow \Spec(A_h)
$$
を得る。$f' \in A_h$ を $f$ の像とする。このとき $(j')^{-1}D(f')$ は $U$ のアフィン開集合
$j^{-1}D(f)$ のうち $g$ が定める主開集合である。したがって $(j')^{-1}D(f)$ はアフィンである。
最後に、$h \in I$ とイデアル $I$ の選び方により
$j'(V) \cap V(f') = j'(j^{-1}D(h) \cap Z)$ は
$\Spec(A_h/(f')) = \Spec((A/f)_h) = D(h) \cap V(f)$ で閉である。
よって補題 \ref{more-groupoids-lemma-get-affine} を適用すると、上で主張したとおり $V$ はアフィンである。
\end{proof}







\section{ind-準アフィン射の降下}
\label{more-groupoids-section-ind-quasi-affine}

\noindent
ind-準アフィン射は『射の続論』節 \ref{more-morphisms-section-ind-quasi-affine} で定義した。
本節は ind-準アフィン射についての『降下』節 \ref{descent-section-quasi-affine} の類似である。

\medskip\noindent
$X$ を準分離スキームとする。$E \subset X$ を、$X$ の準コンパクト開集合の空でない族の
共通部分である部分集合とする。$U_i \subset X$ を準コンパクト開集合、$I$ を空でない集合として
$E = \bigcap_{i \in I} U_i$ と書く。有限交叉を追加することで、任意の $i, j \in I$ に対し
$U_k \subset U_i \cap U_j$ を満たす $k \in I$ が存在すると仮定してよい。この状況で、$X$ 上に
定義された任意の層 $\mathcal{F}$ に対して
\begin{equation}
\label{more-groupoids-equation-sections-of-intersection}
\Gamma(E, \mathcal{F}|_E) = \colim \Gamma(U_i, \mathcal{F}|_{U_i})
\end{equation}
が成り立つ。実際、$i_0 \in I$ を固定し、$X$ を $U_{i_0}$ で、$I$ を
$\{i \in I \mid U_i \subset U_{i_0}\}$ で置き換える。このとき $X$ は準コンパクトかつ準分離で、
したがってスペクトル空間である。『性質』補題
\ref{properties-lemma-quasi-compact-quasi-separated-spectral} を参照せよ。すると等式は『位相』補題
\ref{topology-lemma-make-spectral-space} と『層』補題 \ref{sheaves-lemma-descend-opens} から従う。
（実際、$\mathcal{F}$ がアーベル層なら、この公式は高次コホモロジー群についても成り立つ。
『コホモロジー』補題 \ref{cohomology-lemma-colimit} を参照せよ。）

\begin{lemma}
\label{more-groupoids-lemma-sits-in-functions}
$X$ を ind-準アフィンスキームとする。$E \subset X$ を $X$ の準コンパクト開集合の空でない族の
共通部分とする。$A = \Gamma(E, \mathcal{O}_X|_E)$ および $Y = \Spec(A)$ とおく。
このとき『スキーム』補題 \ref{schemes-lemma-morphism-into-affine} の標準射
$$
j : (E, \mathcal{O}_X|_E) \longrightarrow (Y, \mathcal{O}_Y)
$$
は同型 $(E, \mathcal{O}_X|_E) \to (E', \mathcal{O}_Y|_{E'})$ を定める。ここで
$E' \subset Y$ は準コンパクト開集合の共通部分である。$W \subset E$ が $X$ で開なら、
$j(W)$ は $Y$ で開である。
\end{lemma}

\begin{proof}
$(E, \mathcal{O}_X|_E)$ は局所環付き空間なので、『スキーム』補題
\ref{schemes-lemma-morphism-into-affine} を $A \to \Gamma(E, \mathcal{O}_X|_E)$ に適用できる。
$I \not = \emptyset$、$U_i \subset X$ を準コンパクト開集合として
$E = \bigcap_{i \in I} U_i$ と書く。任意の $i, j \in I$ に対し
$U_k \subset U_i \cap U_j$ を満たす $k \in I$ が存在すると仮定してよい。
$A_i = \Gamma(U_i, \mathcal{O}_{U_i})$ とおく。可換図式
$$
\xymatrix{
(E, \mathcal{O}_X|_E) \ar[r] \ar[d] &
(\Spec(A), \mathcal{O}_{\Spec(A)}) \ar[d] \\
(U_i, \mathcal{O}_{U_i}) \ar[r] &
(\Spec(A_i), \mathcal{O}_{\Spec(A_i)})
}
$$
を得る。$U_i$ は準アフィンなので $U_i \to \Spec(A_i)$ は準コンパクトな開埋め込みである。
一方 $A = \colim A_i$ である。したがって位相空間として
$\Spec(A) = \lim \Spec(A_i)$ である（『極限』補題 \ref{limits-lemma-topology-limit}）。
$E = \lim U_i$ なので（『位相』補題 \ref{topology-lemma-make-spectral-space}）、
$E \to \Spec(A)$ はその像 $E'$ への同相写像であり、$E'$ は $\Spec(A)$ における開集合
$U_i \subset \Spec(A_i)$ の逆像の共通部分である。任意の $e \in E$ に対し、局所環
$\mathcal{O}_{X, e}$ は $\mathcal{O}_{U_i, e}$ の値であり、これは $\Spec(A)$ 上の値と同じである。

\medskip\noindent
補題の最後の主張を示す。$U_i \subset U_j$ を満たす $i, j \in I$ を取る。可換図式
$$
\xymatrix{
U_i \ar[r] \ar[d] & \Spec(A_i) \ar[d] \\
U_i \ar[r] & \Spec(A_j)
}
\quad\quad
\xymatrix{
W \ar[r] \ar[d] & \Spec(A_i) \ar[d] \\
W \ar[r] & \Spec(A_j)
}
\quad\quad
\xymatrix{
W \ar[r] \ar[d] & \Spec(A) \ar[d] \\
W \ar[r] & \Spec(A_j)
}
$$
を考える。『性質』補題 \ref{properties-lemma-cartesian-diagram-quasi-affine} により最初の図式は
Cartesian である。したがって二番目も Cartesian である。極限へ移ると三番目の図式も
Cartesian となるので、この図式の上の水平射は開埋め込みである。
\end{proof}

\begin{lemma}
\label{more-groupoids-lemma-affine-base-change}
スキームの Cartesian 図式
$$
\xymatrix{
X \ar[d]_f \ar[r] & \Spec(B) \ar[d] \\
Y \ar[r] & \Spec(A)
}
$$
が与えられたとする。$E \subset Y$ を $Y$ の準コンパクト開集合の空でない族の共通部分とする。
$Y$ が準分離で $A \to B$ が平坦なら
$$
\Gamma(f^{-1}(E), \mathcal{O}_X|_{f^{-1}(E)}) =
\Gamma(E, \mathcal{O}_Y|_E) \otimes_A B
$$
である。
\end{lemma}

\begin{proof}
$V_i \subset Y$ を準コンパクト開集合として $E = \bigcap_{i \in I} V_i$ と書く。
任意の $i, j \in I$ に対し $V_k \subset V_i \cap V_j$ を満たす $k \in I$ が存在すると
仮定してよい。同様に $X$ で $f^{-1}(E) = \bigcap_{i \in I} f^{-1}(V_i)$ である。
したがって結果は等式 (\ref{more-groupoids-equation-sections-of-intersection}) と、$V_i$ および
$f^{-1}(V_i)$ に対する対応する結果、すなわち『スキームのコホモロジー』補題
\ref{coherent-lemma-flat-base-change-cohomology} から従う。
\end{proof}

\begin{lemma}[Gabber]
\label{more-groupoids-lemma-ind-quasi-affine}
$S$ をスキームとし、$\{X_i \to S\}_{i\in I}$ を fpqc 被覆とする。
$(V_i/X_i, \varphi_{ij})$ を $\{X_i \to S\}$ に関する降下データとする。『降下』定義
\ref{descent-definition-descent-datum-for-family-of-morphisms} を参照せよ。各射
$V_i \to X_i$ が ind-準アフィンならば、この降下データは有効である。
\end{lemma}

\begin{proof}
ind-準アフィンであることは任意の基底変換で保たれるスキームの射の性質である。
『射の続論』補題 \ref{more-morphisms-lemma-base-change-ind-quasi-affine} を参照せよ。
したがって『降下』補題 \ref{descent-lemma-descending-types-morphisms} が適用でき、fpqc 被覆が
アフィン間の一つの平坦全射 $\{X \to S\}$ で与えられる場合に補題を示せば十分である。
$X = \Spec(A)$、$S = \Spec(R)$ と書く。このとき $R \to A$ は忠実平坦な環準同型である。
$(V, \varphi)$ を $S$ 上の $X$ に関する降下データとし、$V \to X$ が ind-準アフィン、
言い換えれば $V$ が ind-準アフィンであると仮定する。

\medskip\noindent
$(U, R, s, t, c)$ を、$U = X$、$R = X \times_S X$ とし、$s$, $t$, $c$ を通常どおりに
取った $S$ 上の亜群スキームとする。『亜群スキーム』補題
\ref{groupoids-lemma-cartesian-equivalent-descent-datum} により、組 $(V, \varphi)$ は亜群スキームの
Cartesian 射 $(U', R', s', t', c') \to (U, R, s, t, c)$ に対応する。
$u' \in U'$ を任意の点とする。『亜群スキーム』補題
\ref{groupoids-lemma-constructing-invariant-opens}、
\ref{groupoids-lemma-first-observation}、\ref{groupoids-lemma-second-observation} により
$u' \in W \subset E \subset U'$ を選べる。ここで $W$ は開かつ $R'$-不変であり、$E$ は
集合論的に $R'$-不変で、準コンパクト開集合の空でない族の共通部分である。

\medskip\noindent
$(V, \varphi)$ の言葉に戻すと、任意の $v \in V$ に対し、次の性質をもつ
$v \in W \subset E \subset V$ が存在する。(a) $W$ は開で
$\varphi(W \times_S X) = X \times_S W$ である。(b) $E$ は準コンパクト開集合の共通部分で、
集合論的に $\varphi(E \times_S X) = X \times_S E$ である。ここで $E \times_S X$ は、射影による
$V \times_S X$ における $E$ の逆像を表し、$X \times_S E$ も同様である。
補題 \ref{more-groupoids-lemma-affine-base-change} により、$\varphi$ は $A \otimes_R A$-代数の同型
\begin{align*}
\Gamma(E, \mathcal{O}_V|_E) \otimes_R A
& =
\Gamma(E \times_S X, \mathcal{O}_{V \times_S X}|_{E \times_S X}) \\
& \to
\Gamma(X \times_S E, \mathcal{O}_{X \times_S V}|_{X \times_S E}) \\
& =
A \otimes_R \Gamma(E, \mathcal{O}_V|_E)
\end{align*}
を定める。これを $\psi$ と呼ぶ。$\varphi$ のコサイクル条件は、『降下』定義
\ref{descent-definition-descent-datum-modules} における $\psi$ のコサイクル条件に移る
（詳細は省略する）。『降下』命題 \ref{descent-proposition-descent-module} により、$R$-代数
$R'$ と $A$-代数の同型
$\chi : R' \otimes_R A \to \Gamma(E, \mathcal{O}_V|_E)$
で、$\psi$ および $R' \otimes_R A$ 上の標準降下データと両立するものが存在する。

\medskip\noindent
補題 \ref{more-groupoids-lemma-sits-in-functions} により、局所環付き空間の標準的な「埋め込み」
$$
j : (E, \mathcal{O}_V|_E) \longrightarrow
\Spec(\Gamma(E, \mathcal{O}_V|_E)) = \Spec(R' \otimes_R A)
$$
を得る。この写像の構成は標準的であり、可換図式
$$
\xymatrix{
& E \times_S X \ar[rr]_\varphi \ar[ld] \ar[rd]^{j'} & &
X \times_S E \ar[rd] \ar[ld]_{j''} \\
E \ar[rd]^j  & &
\Spec(R' \otimes_R A \otimes_R A) \ar[ld] \ar[rd] & &
E \ar[ld]_j \\
& \Spec(R' \otimes_R A) \ar[rd] && \Spec(R' \otimes_R A) \ar[ld] \\
& & \Spec(R')
}
$$
を得る。ここで $j'$, $j''$ は、$\chi$ と $\psi$ の構成に用いた同一視を介し、
$E \times_S X \subset V \times_S X$ と $X \times_S E \subset X \times_S V$ に同じ構成を
適用して得られる。したがって $j(W)$ は $\Spec(R' \otimes_R A)$ の開部分スキームであり、
二つの射影 $\Spec(R' \otimes_R A \otimes_R A) \to \Spec(R' \otimes_R A)$ による逆像は等しい。
『降下』補題 \ref{descent-lemma-open-fpqc-covering} により、$\Spec(A)$ への基底変換が $j(W)$ と
なる開集合 $W_0 \subset \Spec(R')$ が存在する。上の図式を考察すると、降下データ
$(W, \varphi|_{W \times_S X})$ は有効である。『降下』補題
\ref{descent-lemma-effective-for-fpqc-is-local-upstairs} により、元の降下データも有効である。
\end{proof}
